\chapter{Formy kwadratowe i przeszkoda chirurgiczna}
\label{ch:przeszkoda-chirurgiczna}
\index{grupa L@Grupa $L$}

W Rozdziale~\ref{ch:chirurgia-geometryczna-21} operacja na
obramowanej sferze zmieniała rozmaitość, a forma przecięć
na jądrze mapy stopnia jeden pokazywała, dlaczego w
\emph{środkowym wymiarze} nie można usuwać klas po jednej.
Teraz zapiszemy to ograniczenie algebraicznie. Otrzymamy
grupę, w której klasa zerowa oznacza możliwość dokończenia
chirurgii przy odpowiednich założeniach wymiarowych.
Odróżnimy twierdzenia o konkretnych macierzach, które
udowodnimy tutaj, od głębokiego twierdzenia chirurgicznego
Walla, którego geometryczne lematy przywołamy jawnie.

\section{Drogi przecięcia i pierścień grupowy}
\label{sec:pi-formy-22}

Jeśli $X$ jest spójna, połóżmy $\pi=\pi_1(X)$ i
$\Lambda=\Z[\pi]$. Element $\Lambda$ jest skończoną
sumą $\sum_g n_g g$, gdzie $n_g\in\Z$. Mnożenie
przedłuża liniowo mnożenie w $\pi$. W przypadku
$\pi=1$ otrzymujemy zwykły pierścień $\Z$.
Niech $w:\pi\to\{\pm1\}$ będzie charakterem
orientacji celu. Na $\Lambda$ wprowadzamy \emph{inwolucję}
\begin{equation}\label{eq:inwolucja-L-22}
 \overline{\sum_g n_g g}
     =\sum_g n_g w(g)g^{-1}.
\end{equation}
Odwrócenie drogi zamienia $g$ na $g^{-1}$, a
$w(g)$ mówi, czy przeniesienie lokalnej orientacji
wzdłuż tej drogi zmienia znak.

\begin{stwierdzenie}[Dlaczego to inwolucja]
\label{prop:inwolucja-L-22}
Przepis~\eqref{eq:inwolucja-L-22} spełnia
$\overline{ab}=\bar b\bar a$ i
$\overline{\bar a}=a$.
\end{stwierdzenie}
\begin{proof}
Wystarczy sprawdzić elementy grupy, ponieważ obie
strony są addytywne. Ponieważ $w$ jest homomorfizmem,
\[
 \overline{gh}=w(g)w(h)h^{-1}g^{-1}
              =(w(h)h^{-1})(w(g)g^{-1})
              =\bar h\bar g.
\]
Ponadto $w(g^{-1})=w(g)$ i $w(g)^2=1$, skąd
$\overline{\bar g}=g$. Przedłużenie liniowe kończy
rachunek.
\end{proof}

Po podniesieniu sfer do nakrycia uniwersalnego
pojedynczy punkt ich przecięcia wnosi do
\emph{współczynnika przecięcia} znak razy element $g\in\pi$.
Znak zależy od orientacji, a $g$ od obranych dróg
od punktu bazowego do sfer. Dlatego dwa przecięcia
o przeciwnych znakach mogą sumować się do
$g-h\ne0$ w $\Lambda$, choć ich zwykła liczba
przecięcia jest zerem. Zmiana dróg zmienia zapis
macierzy, lecz nie jej klasę izometrii.

\begin{definicja}[Forma $\epsilon$-hermitowska]
\label{def:hermitowska-22}
Niech $P$ będzie skończenie generowanym rzutowym
lewym modułem
nad $\Lambda$ oraz $\epsilon\in\{\pm1\}$.
\emph{Formą $\epsilon$-hermitowską} nazywamy
odwzorowanie $\lambda:P\times P\to\Lambda$
addytywne w każdej zmiennej i takie, że
\[
 \lambda(ax,by)=a\lambda(x,y)\bar b,\qquad
 \lambda(y,x)=\epsilon\overline{\lambda(x,y)}.
\]
Jest \emph{nieosobliwa}, gdy odwzorowanie
$P\to P^*$, $y\mapsto\lambda(-,y)$, jest
izomorfizmem. Tutaj
$P^*=\operatorname{Hom}_\Lambda(P,\Lambda)$
jest lewym modułem z działaniem
$(a\cdot h)(x)=h(x)\bar a$; dzięki temu
$y\mapsto\lambda(-,y)$ jest liniowe.
\end{definicja}

Dla $n=2k$ przecięcia $k$-wymiarowych sfer
spełniają ten wzór z $\epsilon=(-1)^k$:
zamiana kolejności dwóch przestrzeni stycznych
w punkcie przecięcia zamienia bloki po $k$ wektorów,
co daje znak $(-1)^{k^2}=(-1)^k$; odwraca także
drogę, stąd kreska nad współczynnikiem.

\section{Samoprzecięcia i formy kwadratowe}
\label{sec:kwadratowa-22}

Sama $\lambda$ nie zapamiętuje, czy reprezentant
klasy daje się przedstawić jako \emph{obramowana,
zanurzona} sfera. Potrzebna jest poprawka
samoprzecięciowa $\mu$; jej precyzyjną postać
podajemy w języku algebraicznym. Dla
$\epsilon=(-1)^k$ oznaczmy
\[
 Q_\epsilon(\Lambda)=
 \Lambda/\{a-\epsilon\bar a:a\in\Lambda\}.
\]
W mianowniku chodzi o podgrupę addytywną, a nie
zwykle o ideał. Nawias $[a]$ oznacza klasę $a$.

\begin{definicja}[Forma $\epsilon$-kwadratowa]
\label{def:forma-kwadratowa-22}
To para $(\lambda,\mu)$ na $P$, gdzie $\lambda$
jest nieosobliwą formą $\epsilon$-hermitowską, a
$\mu:P\to Q_\epsilon(\Lambda)$ spełnia
\begin{align}
 \mu(x+y)&=\mu(x)+\mu(y)+[\lambda(x,y)],
                   \label{eq:quad-add-22}\\
 \mu(ax)&=[a m\bar a]\quad\text{dla }m\in\mu(x),
                                                \label{eq:quad-scal-22}\\
 \lambda(x,x)&=a+\epsilon\bar a
       \quad\text{dla dowolnego }a\in\mu(x).
                   \label{eq:quad-diag-22}
\end{align}
Drugi wzór jest niezależny od wyboru $m$:
$a(c-\epsilon\bar c)\bar a$ ma postać
$d-\epsilon\bar d$ dla $d=ac\bar a$.
Ostatni wzór jest niezależny od reprezentanta
$a$: dodanie $c-\epsilon\bar c$ nie zmienia
$a+\epsilon\bar a$.
\end{definicja}

Geometrycznie $\mu(x)$ liczy samoprzecięcia
obramowanej immersji sfery i uwzględnia obramowanie.
Równanie~\eqref{eq:quad-add-22} mówi:
gdy połączymy reprezentantów $x$ i $y$, nowe
samoprzecięcia obejmują też przecięcia obu sfer.
W \eqref{eq:quad-diag-22} widać, dlaczego
„połowa przekątnej” nie zawsze wystarcza:
nad $\Z[\pi]$ nie można po prostu dzielić przez $2$.

\begin{stwierdzenie}[Rachunek przecięć na nakryciu]
\label{prop:rachunek-przeciec-Wall-22}
Niech $k\geq3$, a $M^{2k}\to X^{2k}$ będzie
$k$-spójnym odwzorowaniem normalnym stopnia jeden.
Wybierzmy obramowane, poprzeczne immersje
$k$-sfer reprezentujące $x,y\in K_k(M)$,
wraz z drogami od punktu bazowego. Każdy punkt
$p$ przecięcia wnosi znak $\varepsilon_p$ i element
$g_p\in\pi_1(X)$, wyznaczony przez przejście
między wybranymi podniesieniami sfer. Wtedy
\[
 \lambda(x,y)=\sum_{p\in x\pitchfork y}
                  \varepsilon_p g_p\in\Lambda.
\]
Wybierając przy każdym podwójnym punkcie immersji
$x$ jedną z dwóch kolejności gałęzi, otrzymujemy
\[
 \mu(x)=\left[\sum_{p\in\operatorname{Dbl}(x)}
                 \varepsilon_p g_p\right]
                 \in Q_{(-1)^k}(\Lambda).
\]
Wartości te nie zależą od regularnej homotopii
reprezentantów; spełniają
\eqref{eq:quad-add-22} i~\eqref{eq:quad-diag-22}.
\end{stwierdzenie}
\begin{proof}
Podnieśmy obie sfery do nakrycia uniwersalnego.
Współczynnik przy $g$ w $\lambda(x,y)$ jest
zwykłą algebraiczną liczbą przecięć jednego
podniesienia z przesunięciem drugiego przez $g$.
Jest więc stały przy homotopii. Po zamianie
obu sfer droga odwraca się, a znak zmienia się
o $(-1)^k w(g)$. To daje hermitowskość z
Definicji~\ref{def:hermitowska-22}.

Przy podwójnym punkcie zamiana kolejności
gałęzi zmienia $\varepsilon_p g_p$ na
$(-1)^k\overline{\varepsilon_p g_p}$.
Oba wyrazy mają tę samą klasę w
$Q_{(-1)^k}(\Lambda)$. Podczas ogólnej regularnej
homotopii punkty podwójne rodzą się lub znikają
parami o przeciwnych znakach i tej samej
etykiecie, zatem $\mu$ się nie zmienia.
Po połączeniu dwóch sfer rurką nowe punkty
podwójne to, oprócz dawnych, dokładnie ich
wzajemne przecięcia. Dostajemy stąd
\eqref{eq:quad-add-22}.

Na koniec odsuńmy immersję reprezentującą $x$
wzdłuż jej obramowania. Przy każdym podwójnym
punkcie odsunięta sfera przecina pierwotną
dwukrotnie, raz dla każdej kolejności gałęzi.
Ich suma wynosi $a+(-1)^k\bar a$ dla składnika
$a$ w $\mu(x)$. Wkład przekątnej z klasy Eulera
wiązki normalnej znika, bo sfera jest obramowana.
To jest~\eqref{eq:quad-diag-22}.
\end{proof}

Rachunek punktów wyjaśnia wzory, lecz aby
traktować $\mu$ jako funkcję na $K_k(M)$,
trzeba jeszcze wiedzieć, że dane normalne
wyznaczają właściwą klasę regularnej homotopii
obramowanych immersji. To identyfikacja
użyta w twierdzeniu 5.2 Walla; nie wynika
z samego zliczenia punktów.

Na przykład dwa przecięcia o znakach $+1$ i $-1$
oraz etykietach $1$ i $t$ dają
$\lambda(x,y)=1-t\in\Z[t,t^{-1}]$ dla
$\pi_1(X)=\langle t\rangle$. Zwykła liczba
przecięcia jest zerem, lecz okrąg złożony z
łuków obu sfer reprezentuje $t$ lub $t^{-1}$
i nie ogranicza
dysku Whitneya. Dopiero przeciwne znaki z
\emph{jednakową} etykietą pozwalają szukać
takiego dysku. Warunek $\mu(x)=0$ dopuszcza
analogiczne sparowanie samoprzecięć. Wyjaśnijmy
osobno, dlaczego w wymiarach co najmniej sześć
takie pary można usunąć.

\subsection*{Od samoprzecięć do zanurzonej sfery}
\index{trik Whitneya@Trik Whitneya!samoprzecięcia}
W tym kroku potrzebna jest \emph{obramowana
immersja}, a nie jedynie ciągłe odwzorowanie
$S^k\to M$. Obramowanie jest częścią klasy
regularnej homotopii wyznaczonej przez mapę
normalną i wypełnienie obrazu sfery w $X$.

\begin{stwierdzenie}[Kryterium zanurzania Walla w środku wymiaru]
\label{prop:osadzenie-Wall-52-22}
Niech $k\geq3$, a $M^{2k}$ będzie spójną gładką
rozmaitością. Niech $u:S^k\times D^k\looparrowright
\operatorname{int}M$ będzie obramowaną immersją,
której rdzeń ma jedynie poprzeczne punkty podwójne.
Oznaczmy $\pi=\pi_1(M)$ i wybierzmy drogi bazowe
potrzebne do określenia
$\mu(u)\in Q_{(-1)^k}(\Z[\pi])$.
Jeśli $\mu(u)=0$, to $u$ jest regularnie homotopijna
do immersji, której rdzeń jest zanurzony. Po
zwężeniu włókien $D^k$ otrzymujemy zanurzenie
$S^k\times D^k\hookrightarrow M$ w tej samej
klasie regularnej homotopii. Dla zanurzonego rdzenia
warunek $\mu(u)=0$ jest konieczny.
\end{stwierdzenie}
\begin{proof}[Dowód z obramowanym trikiem Whitneya]
\textbf{1. Parowanie punktów.}
Niech $\epsilon=(-1)^k$. W grupie
$Q_\epsilon(\Z[\pi])$ zamiana kolejności dwóch
gałęzi w punkcie o etykiecie $g$ zamienia składnik
$\varepsilon g$ na $\epsilon\overline{\varepsilon g}
=\epsilon w(g)\varepsilon g^{-1}$.
Rozłóżmy bazę grupy pierścienia na orbity
$g\leftrightarrow g^{-1}$. Dla orbity dwuelementowej
możemy więc wszystkie składniki zapisać przy tej
samej etykiecie $g$; ich suma całkowita jest zerowa.
Dla orbity jednoelementowej $g=g^{-1}$ otrzymujemy
albo kopię $\Z$, albo $\Z/2$, zależnie od znaku
$\epsilon w(g)$. W pierwszym przypadku suma znaków
jest zerowa. W drugim liczba punktów jest parzysta,
a zamiana gałęzi zmienia ich znaki. W każdym
przypadku punkty dają się zatem podzielić na pary
$(p,q)$ o \emph{jednakowej} etykiecie i przeciwnych
znakach, po odpowiednim wyborze kolejności gałęzi.

\textbf{2. Okrąg Whitneya.}
W dziedzinie $S^k$ wybierzmy od $p$ do $q$ jeden
łuk na pierwszym płacie, a drugi na drugim.
Ponieważ $k\geq3$, łuki można wybrać rozłączne
poza końcami i ominąć preobrazy pozostałych
punktów podwójnych. Ich obrazy tworzą okrąg
Whitneya. Przy ustalonych drogach bazowych jego
klasa w $\pi_1(M)$ jest, z dokładnością do
odwrócenia i sprzężenia, iloczynem $g_p g_q^{-1}$.
Równość etykiet daje więc pętlę homotopijnie
trywialną. Ten krok nie wymaga $\pi_1(M)=1$;
właśnie temu służą etykiety w $\Z[\pi]$.

\textbf{3. Czysty dysk.}
Odsuńmy okrąg od obu lokalnych płatów. Inkluzja
$M\setminus u(S^k\times\{0\})\hookrightarrow M$
indukuje izomorfizm na $\pi_1$: pętle i homotopie
reprezentowane przez dwuwymiarowe dyski można
przesunąć poza obraz sfery, gdyż
$2+k<2k$ dla $k\geq3$. Skończony zbiór punktów
podwójnych nie zmienia tego rachunku. Odsunięty
okrąg ogranicza więc dysk w dopełnieniu.
Transwersalność względna pozwala wybrać dysk
gładki, z wnętrzem rozłącznym z obrazem rdzenia;
nierówność $2+2<2k$ pozwala usunąć także jego
samoprzecięcia. Dołączając kołnierz przy łukach,
otrzymujemy zanurzony dysk Whitneya $D_W$ o brzegu
na dwóch wybranych płatach.

\textbf{4. Obramowanie i ruch.}
Sam czysty dysk nie wystarcza do ruchu z
Stwierdzenia~\ref{prop:lokalny-Whitney-21}.
Wzdłuż pierwszego łuku wybieramy $(k-1)$ kierunków
stycznych do płata i prostopadłych do łuku;
w punktach końcowych poprzeczność identyfikuje
je z kierunkami normalnymi do drugiego płata.
Ponieważ okrąg Whitneya jest zerowy w $\pi_1(M)$,
możemy spójnie przenosić orientację $M$ wzdłuż
jego brzegu także wtedy, gdy całe $M$ nie jest
orientowalne. Przeciwne znaki w $p$ i $q$ sprawiają,
że po odwróceniu kierunku drugiego łuku orientacje
ram końcowych są zgodne. Przestrzeń pośrednich
zorientowanych ram jest spójna, bo $SO(k-1)$
jest spójna dla $k\geq3$; wybieramy więc ramę
normalną do drugiego płata na całym drugim łuku.
Obie części sklejają się w $(k-1)$-ramę na
$\partial D_W$. Wiązka normalna
dysku ma rangę $2k-2$ i jest trywialna. Przeszkodę
przedłużenia powyższej $(k-1)$-ramy mierzy
$\pi_1(V_{2k-2,k-1})$, która jest zerowa dla
$k\geq3$. Pozostałe kierunki uzupełniamy nad
dyskiem. Otrzymujemy lokalne otoczenie modelowe
$D^2\times D^{k-1}\times D^{k-1}$, w którym
dwa płaty zajmują dwa prostopadłe układy pasów.
Precyzyjna konstrukcja tego otoczenia jest
obramowanym trikiem Whitneya: używamy
\href{https://webhomes.maths.ed.ac.uk/~v1ranick/surgery/hcobord.pdf}
{Milnora, twierdzenia 6.6 i lematów 6.7 oraz 6.13}.
W naszym zastosowaniu oba płaty mają wymiar $k$
w otoczeniu wymiaru $2k$. Warunki wymiarowe
Milnora sprowadzają się więc do $k\geq3$.
Osobno potrzebujemy ściągalności okręgu Whitneya;
jego łuki wybieramy tak, by omijały pozostałe
przecięcia. Lematy 6.7
i 6.13 dostarczają właśnie produktowego modelu
otoczenia dysku oraz zgodnych ram normalnych.
Twierdzenie 6.6 Milnora jest sformułowane dla
dwóch zanurzonych podrozmaitości; zastosowanie
modelu do dwóch lokalnych płatów jednej immersji
wymaga jeszcze argumentu Walla z twierdzenia 5.2.
Przed tym krokiem można zwęzić włókna $D^k$,
aby otoczenie dysku było wolne od pozostałych
części pogrubionej immersji.
Lokalny ruch z
Stwierdzenia~\ref{prop:lokalny-Whitney-21}
usuwa teraz $p$ i $q$. Dla jednej samoprzecinającej
się sfery ograniczamy uwagę do dwóch zanurzonych
fragmentów odpowiadających wybranym łukom
i wykonujemy go na jednym \emph{lokalnym płacie};
jest to regularna homotopia
immersji, nie izotopia otaczająca całego $M$.
Jej podparcie omija pozostałe punkty podwójne.
Przejście od twierdzenia Milnora o dwóch zanurzonych
podrozmaitościach do dwóch płatów jednej immersji
jest geometryczną częścią kryterium~5.2 Walla.
Powtarzamy krok dla wszystkich par.

Regularna homotopia przenosi obramowanie i klasę
normalną. Gdy rdzeń jest już zanurzony, zwężamy
włókna $D^k$ w obramowanym otoczeniu; zwartość
$S^k$ zapewnia, że dla dostatecznie małego promienia
całe pogrubienie jest zanurzone. Zanurzony rdzeń
nie ma punktów podwójnych, więc jego $\mu$ wynosi
zero.
\end{proof}

Przy zastosowaniu do $x\in K_k(M)$ potrzebne są
jeszcze dwa wejścia. Po przygotowaniu
$k$-spójnej mapy normalnej $M\to X$ mamy
$\pi_1(M)\cong\pi_1(X)$ oraz
$K_k(M)\cong\pi_{k+1}(f)$; wybór właściwej
obramowanej klasy regularnej homotopii immersji
$S^k\times D^k$ jest treścią
\href{https://webhomes.maths.ed.ac.uk/~v1ranick/books/scm.pdf}
{Walla, twierdzenia 1.1}. Przytoczona konstrukcja
produktowego otoczenia dysku pozostaje osobnym
lematem geometrycznym Milnora; powyższy rachunek
nie jest samodzielnym dowodem tego lematu.
W wymiarze cztery nierówności z kroku 3 zawodzą,
więc sam warunek $\mu=0$ nie daje analogicznego
wniosku. Ponadto uogólnienie $\lambda,\mu$ z
uwagi 5.2.1 Walla na nakrycie indukowane z innej
rozmaitości nie obejmuje kryterium zanurzania.

\begin{przyklad}[Co zostaje, gdy $\pi=1$]
\label{prz:mu-Z-22}
Dla $\Lambda=\Z$ inwolucja jest zwykłą tożsamością.
Jeśli $k$ jest parzyste, $Q_+(\Z)=\Z$,
a~\eqref{eq:quad-diag-22} daje
$\lambda(x,x)=2\mu(x)$. Forma jest więc
\emph{parzysta}; $\mu(x)=\lambda(x,x)/2$.
Jeśli $k$ jest nieparzyste,
$Q_-(\Z)=\Z/2$; forma $\lambda$ jest
antysymetryczna i $\lambda(x,x)=0$,
ale $\mu(x)\in\Z/2$ może być niezerowe.
Ta ostatnia informacja nie daje się odczytać
z samej macierzy przecięć.
\end{przyklad}

\begin{definicja}[Para hiperboliczna i lagranżjan]
\index{para hiperboliczna i lagranzjan@Para hiperboliczna i lagranżjan}
\label{def:hiperboliczna-22}
Wolny moduł o bazie $(e_1,f_1,\dots,e_r,f_r)$
z macierzą w każdej parze
\[
 H_\epsilon=
 \begin{pmatrix}0&1\\\epsilon&0\end{pmatrix},
 \qquad \mu(e_i)=\mu(f_i)=0,
\]
oraz z zerowymi sprzężeniami między różnymi
parami, nazywamy \emph{hiperbolicznym}.
\emph{Lagranżjanem} formy kwadratowej jest
bezpośredni składnik $L\subset P$ taki, że
$L=L^\perp$ oraz $\mu|_L=0$.
W parze hiperbolicznej
$L=\langle e_1,\dots,e_r\rangle$ jest lagranżjanem.
\end{definicja}

Parę o macierzy $H_\epsilon$ realizują klasy
dwóch sfer czynnikowych w $S^k\times S^k$,
jak w Przykładzie~\ref{prz:forma-produkt-21}.
Lokalny model ich jednego przecięcia
pokazuje Rysunek~\ref{fig:para-hiperboliczna-22}.

\begin{figure}[htbp]
\centering
\begin{tikzpicture}[>=Latex,font=\small,line cap=round]
 \fill[manifoldblue!12,rounded corners=4pt]
   (-3.15,-.20) rectangle (3.15,.20);
 \fill[manifoldgold!19,rounded corners=4pt]
   (-.20,-1.12) rectangle (.20,1.12);
 \draw[manifoldblue!80!black,very thick]
   (-3.15,0)--(3.15,0);
 \draw[manifoldgold!85!black,very thick]
   (0,-1.12)--(0,1.12);
 \fill[manifoldred!85!black] (0,0) circle (2.5pt);
 \node[manifoldblue!80!black,anchor=west]
   at (-3.15,-.65) {$D^k\times\{0\}$};
 \node[manifoldgold!75!black,anchor=west]
   at (.38,.88) {$\{0\}\times D^k$};
 \node[manifoldred!85!black,anchor=west]
   at (.32,-.43) {jedno przecięcie};
\end{tikzpicture}
\caption{Lokalny model sfer
$e=S^k\times\{*\}$ i $f=\{*\}\times S^k$.
Ich dyski współrzędnych przecinają się
tylko w środku; po uzgodnieniu orientacji
$\lambda(e,f)=1$. Każda sfera ma
trywialną wiązkę normalną, więc
$\mu(e)=\mu(f)=0$. Jasne pasy pokazują
jedynie przekrój otoczeń, a nie całe sfery.}
\label{fig:para-hiperboliczna-22}
\end{figure}

\subsection*{Przesunięcie uchwytu jako zmiana bazy}
Stwierdzenie~\ref{prop:slide-18} opisało geometryczny
ruch dwóch uchwytów tego samego indeksu:
$e_i$ zostaje zastąpione przez $e_i\pm e_j$.
W wymiarze środkowym wpływ tego ruchu na
rachunek przecięć można wypisać bez kolejnego
twierdzenia geometrycznego. Jeżeli kolumny
macierzy $P$ wyrażają nową bazę w starej,
to macierz formy w nowej bazie wynosi
\begin{equation}\label{eq:slide-forma-22}
 A'=P^{T}AP,\qquad
 \mu(e_i\pm e_j)
   =\mu(e_i)+\mu(e_j)
      \pm[\lambda(e_i,e_j)].
\end{equation}
Pierwsza równość pochodzi z dwuliniowości
$\lambda$ nad $\Z$; druga z
\eqref{eq:quad-add-22} i
$\mu(-e_j)=\mu(e_j)$.
Na przykład dla symetrycznej pary hiperbolicznej
$A=\left(\begin{smallmatrix}0&1\\1&0\end{smallmatrix}\right)$
przesunięcie $e'=e+f$ daje
\[
 A'=\begin{pmatrix}2&1\\1&0\end{pmatrix},
 \qquad \mu(e')=1.
\]
Pojawienie się niezerowego wyrazu na przekątnej
nie tworzy nowej przeszkody: zmieniła się baza
tej samej formy. Przy niebanalnej $\pi_1$ w
\eqref{eq:slide-forma-22} trzeba zastąpić
transpozycję transpozycją sprzężoną oraz
śledzić etykiety dróg; sama macierz całkowita
nie zawiera wtedy całej informacji.

Wypiszmy taką zmianę bazy nad
$\Lambda=\Z[\pi]$ z inwolucją
\eqref{eq:inwolucja-L-22}. W sumie dwóch
płaszczyzn hiperbolicznych połóżmy, dla
dowolnego $a\in\Lambda$,
\[
 e'_1=e_1+a e_2,\qquad e'_2=e_2,\qquad
 f'_1=f_1,\qquad f'_2=f_2-\bar a f_1.
\]
Współczynnik $\bar a$ jest konieczny:
seskwiliniowość daje
\[
 \lambda(e'_1,f'_2)
 =\lambda(ae_2,f_2)-\lambda(e_1,\bar a f_1)
 =a-\overline{\bar a}=0.
\]
Pozostałe mieszane parowania nadal wynoszą
$\lambda(e'_i,f'_j)=\delta_{ij}$; parowania
wewnątrz rodzin $e'$ i $f'$ są zerowe.
Z~\eqref{eq:quad-add-22} wynika również
$\mu(e'_i)=\mu(f'_i)=0$, bo składniki dodawane
do siebie należą do wzajemnie ortogonalnych
lagranżjanów. Zatem nowa baza jest znowu
hiperboliczna. Na $\langle e_1,e_2\rangle$
zmianę opisuje macierz elementarna z
współczynnikiem $a$, a na dualnej rodzinie
odwrotna zmiana sprzężona ze współczynnikiem
$-\bar a$. Obie mają klasę zerową w
$\operatorname{Wh}(\pi)$
(Definicja~\ref{def:k1-wh-20}). Ruch nie zmienia
więc \emph{prostej} klasy bazy; sama równość
macierzy form bez tej kontroli nie dawałaby
takiego wniosku.

\begin{stwierdzenie}[Dlaczego hiperboliczna para nie przeszkadza]
\label{prop:hiperboliczna-22}
W parze $(e,f)$ mamy $\lambda(e,f)=1$,
$\lambda(e,e)=\lambda(f,f)=0$ oraz
$\mu(e)=\mu(f)=0$. Formy hiperboliczne
tworzą neutralne składniki przy stabilizacji.
Ich sygnatura, gdy $\epsilon=+1$ i $\Lambda=\Z$,
jest równa zeru.
\end{stwierdzenie}
\begin{proof}
Pierwsze równości odczytujemy z macierzy.
Na prostą generowaną przez $e$ forma
i funkcja kwadratowa znikają, a jej
ortogonalne dopełnienie wyznacza warunek
$\lambda(ae+bf,e)=\epsilon b=0$, więc
to znowu $\langle e\rangle$.
Dla $\epsilon=+1$ macierz ma wartości własne
$+1$ i $-1$, stąd jej sygnatura wynosi zero.
Suma wielu par ma lagranżjan będący sumą
odpowiednich prostych.
\end{proof}

\begin{figure}[htbp]
\centering
\begin{tikzpicture}[>=Latex,font=\small,line cap=round]
 \shade[inner color=white,outer color=manifoldblue!18,
        rounded corners=8pt] (-2.7,-1.08) rectangle (2.7,1.08);
 \draw[manifoldblue!60!black,rounded corners=8pt]
       (-2.7,-1.08) rectangle (2.7,1.08);
 \draw[manifoldblue!70!black,densely dashed]
       (-2.30,.72)--(2.30,.72);
 \draw[manifoldred!82!black,very thick,->]
       (-2.06,0)--(2.10,0)
       node[near end,below] {$e$};
 \draw[manifoldgreen!75!black,very thick,->]
       (0,-.81)--(0,.80)
       node[near end,right] {$f$};
 \fill[manifoldgold!85!black] (0,0) circle (2.9pt);
 \node[anchor=west] at (3.03,.48)
       {$\lambda(e,f)=1$};
 \node[anchor=west] at (3.03,-.03)
       {$\mu(e)=\mu(f)=0$};
 \node[anchor=west] at (3.03,-.55)
       {$\langle e\rangle=\langle e\rangle^\perp$};
\end{tikzpicture}
\caption{Lokalny schemat pary hiperbolicznej.
Sfery reprezentujące $e$ i $f$ spotykają się poprzecznie
w jednym punkcie, a każdą z nich osobno można
obramować bez samoprzecięć. Wybór kierunku $e$
wyznacza lagranżjan; prostokąt jest tylko fragmentem
otoczenia przecięcia, nie całą rozmaitością.}
\label{fig:hiperboliczna-22}
\end{figure}

\begin{przyklad}[Produkt sfer z poprzedniego rozdziału]
\label{prz:produkt-L-22}
W Przykładzie~\ref{prz:forma-produkt-21}
klasy $e=[S^k\times\{*\}]$ i
$f=[\{*\}\times S^k]$ mają macierz $H_{(-1)^k}$.
Dla standardowych danych normalnych lokalnej
stabilizacji z
Lematu~\ref{lem:wall-lokalna-stabilizacja-22}
obie sfery są obramowanymi reprezentantami
bez punktów podwójnych, więc
$\mu(e)=\mu(f)=0$. Samą macierz przecięć
trzeba uzupełnić o ten wybór obramowań;
Przykład~\ref{prz:Arf-dwie-22} pokazuje dlaczego.
To przykład jądra \emph{niezerowego jako moduł},
które daje \emph{zerową przeszkodę}:
można usunąć parę za pomocą odpowiednich chirurgii.
Dla $k=3$ jest to jądro mapy na górnym brzegu
śladu z Przykładu~\ref{prz:slad-s6-produkt-21}.
Jego miejsce w ciągu chirurgii opisuje
Przykład~\ref{prz:slad-s6-ciag-23}.
\end{przyklad}

\section{Co oznacza grupa \texorpdfstring{$L$}{L}}
\label{sec:grupy-L-22}

W wymiarze $2k$ sumujemy formy kwadratowe
ortogonalnie. Dwie formy uważa się za równoważne
w sensie Witta, jeśli ich ortogonalna różnica,
po ewentualnym dodaniu form hiperbolicznych,
ma lagranżjan. W wersji $L^s$ moduły mają
preferowane bazy stabilne, a dla lagranżjanu
$L\subset P$ odwzorowanie $P/L\to L^*$
wyznaczone przez $\lambda$ ma być \emph{prostym}
izomorfizmem względem tych baz. Jest to algebraiczny warunek
zerowania różnicy klas. Klasy z dodawaniem
przez sumę ortogonalną tworzą \emph{kwadratową grupę
chirurgiczną} $L_{2k}(\Lambda)$, z odpowiednią
dekoracją opisaną poniżej. Wersję opartą na bazach
i prostych równoważnościach oznaczymy
$L^s_{2k}(\Lambda,w)$. Warunek „prosta” dotyczy
również dualności formy: izomorfizm do modułu dualnego
ma zerową torsję w $\operatorname{Wh}(\pi)$.
Wersja $L^h_{2k}$ używa zwykłej równoważności
homotopijnej. Dla $\pi=1$ te rozróżnienia znikają.

\subsection*{Dekoracje $s$ i $h$ w rachunku baz}
\label{sec:dekoracje-L-22}
Intuicją dekoracji jest pytanie, jaki rodzaj
równoważności ma pozostać po usunięciu jądra:
$h$ oznacza równoważność homotopijną, a $s$
równoważność \emph{prostą}. W rozdziale o torsji
warunek prostoty zapisaliśmy jako znikanie klasy
macierzy w $\operatorname{Wh}(\pi)$
(Definicja~\ref{def:k1-wh-20}). Dla formy
na bazowanym wolnym module jej dualność
$P\to P^*$ musi być w wersji $s$ prostym
izomorfizmem. Przy usuwaniu lagranżjanu
trzeba ponadto sprawdzić prostotę
$P/L\to L^*$, co wykorzystaliśmy w
Lemacie~\ref{lem:lagr-hip-Z-22}.
W wersji $h$ wymagamy zwykłych izomorfizmów:
macierz odwracalna może mieć niezerową klasę
Whiteheada. To rozróżnienie i odpowiadające
mu grupy przeszkód opisuje
\href{https://webhomes.maths.ed.ac.uk/~v1ranick/books/scm.pdf}
{Wall, \emph{Surgery on Compact Manifolds}, §17D}.

\begin{przyklad}[Odwracalna dualność, która nie jest prosta]
\label{prz:dekoracje-C5-22}
Niech $\pi=C_5=\langle t\mid t^5=1\rangle$,
$w=1$, $\Lambda=\Z[C_5]$ oraz
$u=1-t-t^{-1}$. W
Przykładzie~\ref{prz:wh-c5-20} sprawdziliśmy,
że $u$ jest jednostką z odwrotnością
$v=1-t^2-t^{-2}$, ale $[u]\ne0$ w
$\operatorname{Wh}(C_5)$. Ponadto $\bar u=u$
i $\bar v=v$. Na $P=\Lambda e\oplus\Lambda f$
weźmy formę $+$-kwadratową określoną przez
\[
 \lambda(e,e)=\lambda(f,f)=0,\qquad
 \lambda(e,f)=\lambda(f,e)=u,\qquad
 \mu(e)=\mu(f)=0.
\]
Dla $x=ae+bf$ mamy
$\mu(x)=[a u\bar b]\in Q_+(\Lambda)$;
stąd $\lambda(x,x)$ jest sumą tego
reprezentanta i jego sprzężenia, zgodnie z
\eqref{eq:quad-diag-22}. Podstawienie
$f'=vf$ daje $\lambda(e,f')=u\bar v=1$
i $\mu(f')=0$, więc po zwykłej zmianie bazy
forma jest hiperboliczna.

W bazie $(e,f)$ macierz dualności ma postać
$\left(\begin{smallmatrix}0&u\\u&0\end{smallmatrix}\right)$.
Zamiana dwóch kolumn wnosi tylko znak,
który znika w grupie Whiteheada. Dlatego
torsja tej dualności wynosi
$[u]+[u]=2[u]\ne0$.
Ostatnią nierówność można sprawdzić tym
samym homomorfizmem co w
Przykładzie~\ref{prz:wh-c5-20}:
$|u(e^{2\pi i/5})|^2\ne1$.
Forma dopuszcza więc zwykłą dualność i
reprezentuje zero w opisie $h$, lecz z tą
preferowaną bazą nie spełnia warunku prostej
dualności potrzebnego do zdefiniowania jej
klasy w $L^s_{4j}(\Lambda,1)$.
Przykład dotyczy danych algebraicznych;
nie twierdzimy tu, że wybrana forma jest
jądrem określonej mapy normalnej.
\end{przyklad}

Gdy $\pi=1$, mamy
$\operatorname{Wh}(\pi)=0$
(Stwierdzenie~\ref{prop:wh-trywialna-20}),
więc powyższa różnica warunków znika.
Relatywna wersja twierdzenia o przeszkodzie
zachowuje to samo rozróżnienie na brzegu:
przy dekoracji $s$ zastana mapa brzegu jest
prostą równoważnością homotopijną, przy
dekoracji $h$ zwykłą równoważnością.

W nieparzystym wymiarze jądro nie jest opisywane
pojedynczą formą środkową. Model algebraiczny ma
postać \emph{formacji}: hiperboliczny moduł
kwadratowy wraz z dwoma lagranżjanami. Ich
przecięcie i iloraz przez ich sumę mierzą dwa
sąsiednie jądra mapy; wyprowadzimy to niżej.
Formacje dzielimy przez stabilizację i relację
wyznaczoną przez kobordyzm formacji; otrzymujemy
$L^s_{2k+1}(\Lambda,w)$ lub $L^h_{2k+1}(\Lambda,w)$.
W pełnej definicji obejmującej wszystkie wymiary
są to klasy kobordyzmu skończonych
\emph{kwadratowych kompleksów Poincarégo}
nad $\Lambda$. Konstrukcja formacji jest ich
nieparzystowym opisem po usunięciu homotopijnych
składników poniżej środka. Nie utożsamiamy
nieparzystej grupy z grupą Witta macierzy z poprzedniej sekcji.

\begin{uwaga}[Nie mylmy dwóch słów „torsja”]
W $\operatorname{Wh}(\pi)$ torsja oznacza torsję
\emph{równoważności łańcuchowej} z
Rozdziału~\ref{ch:s-kobordyzm}; w grupie $L$
klasyfikujemy formy i kwadratowe kompleksy.
To, że algebraiczne $K$- i $L$-teoria korzystają
z tych samych modułów nad $\Z[\pi]$, nie czyni
ich przeszkód tym samym obiektem.
\end{uwaga}

\section{Przypadek prostej spójności: sygnatura i Arf}
\label{sec:obliczenia-LZ-22}

Dla $\pi=1$ grupy kwadratowe w wysokowymiarowej
teorii chirurgii mają czterookresowy wzór
\begin{equation}\label{eq:L-Z-22}
 L_{4j}(\Z)\cong\Z,\qquad
 L_{4j+1}(\Z)=0,\qquad
 L_{4j+2}(\Z)\cong\Z/2,\qquad
 L_{4j+3}(\Z)=0.
\end{equation}
Wzór jest twierdzeniem klasyfikacyjnym teorii $L$;
poniżej wyjaśniamy jego dwa niezerowe inwarianty
i sprawdzamy konkretne formy. Nie wyprowadzamy
samych obliczeń grup nieparzystych z rachunku
dwóch macierzy: wymagałoby to teorii formacji.

\subsection*{Wymiary podzielne przez cztery}
Niech $n=4j$, $k=2j$. Środkowa forma jest
parzysta, symetryczna i unimodularna.
Jej \emph{sygnatura} to liczba dodatnich minus
liczba ujemnych wartości własnych nad $\R$.
Klasyczny rachunek form całkowitych daje
\[
 \sigma(f,b)=\frac{\operatorname{sign}\lambda_K}{8}
       \in L_{4j}(\Z)\cong\Z.
\]
Podzielność sygnatury przez $8$ dla parzystej
formy unimodularnej oraz to, że powyższy wzór
daje izomorfizm grup Witta, to nietrywialne
twierdzenia o formach kwadratowych. Nie wynikają
z samego zdiagonalizowania formy nad $\R$.
Sygnatura hiperbolicznej pary jest zerowa, jak
wykazaliśmy w Stwierdzeniu~\ref{prop:hiperboliczna-22}.

\begin{przyklad}[Forma $E_8$ ma niezerową klasę]
\label{prz:E8-L-22}
Niech $E_8$ będzie macierzą symetryczną $8\times8$:
na przekątnej ma $2$, przy krawędziach
\[
 (1,3),(3,4),(4,5),(5,6),(6,7),(7,8),(2,4)
\]
ma $-1$, a poza nimi zera. Schemat jej niezerowych
wyrazów pozadiagonalnych ma postać
\[
\begin{tikzpicture}[x=.72cm,y=.65cm,baseline=(current bounding box.center)]
 \draw[manifoldblue!68!black,thick]
   (0,0)--(1,0)--(2,0)--(3,0)--(4,0)--(5,0)--(6,0)
   (2,0)--(2,1);
 \foreach \num/\x/\y in {1/0/0,3/1/0,4/2/0,5/3/0,
                            6/4/0,7/5/0,8/6/0,2/2/1}
   \node[circle,draw=manifoldblue!75!black,fill=white,
         inner sep=1.5pt,font=\scriptsize] at (\x,\y) {$\num$};
\end{tikzpicture}
\]
Jej wiodące minory główne wynoszą kolejno
$2,4,6,5,4,3,2,1$. Można to sprawdzić, kolejno
rozwijając wyznacznik po liściu narysowanego
drzewa (Definicja~\ref{def:drzewo-grafowe}):
jeśli usuwamy liść $v$ o jedynym
sąsiedzie $u$, wyznacznik jest
$2\det A_{G-v}-\det A_{G-\{u,v\}}$.
Zatem kryterium Sylvestera daje dodatnią
określoność, a ostatni minor $1$ daje
unimodularność. Dla całkowitego wektora $x$
liczba $x^TE_8x$ jest parzysta, bo składniki
pozadiagonalne występują podwojone. Definiujemy
$\mu(x)=x^TE_8x/2$. Otrzymujemy formę
kwadratową o sygnaturze $8$, a więc
$[E_8]=1\in L_{4j}(\Z)$.
Jest to \emph{przykład algebraiczny}; realizowanie
zadanej formy przez konkretną gładką mapę
normalną wymaga osobnego twierdzenia realizacyjnego.
\end{przyklad}

\subsection*{Wymiary $4j+2$}
Niech $k=2j+1$. Redukcja $P$ modulo $2$
daje przestrzeń symplektyczną
$V=P/2P$ nad $\mathbb F_2$ z formą
$b=\lambda\bmod2$. Funkcja $q=\mu\bmod2$
spełnia
\begin{equation}\label{eq:Arf-q-22}
 q(x+y)=q(x)+q(y)+b(x,y).
\end{equation}
Dla bazy symplektycznej $(e_1,f_1,\dots,e_r,f_r)$
definiujemy
\[
 \operatorname{Arf}(q)
       =\sum_{i=1}^r q(e_i)q(f_i)\pmod2.
\]

\begin{stwierdzenie}[Niezależność od bazy]
\label{prop:Arf-baza-22}
Podany wzór na $\operatorname{Arf}(q)$ nie zależy
od wyboru bazy symplektycznej. Suma ortogonalna
dodaje inwarianty Arfa.
\end{stwierdzenie}
\begin{proof}
Policzmy liczbę niezależną od bazy:
\[
 S(q)=\sum_{x\in V}(-1)^{q(x)}.
\]
Z~\eqref{eq:Arf-q-22} i rozkładu na ortogonalne
płaszczyzny $\langle e_i,f_i\rangle$ wynika,
że $q$ na jednej płaszczyźnie przyjmuje na
$ae_i+bf_i$ wartość
$a q(e_i)+bq(f_i)+ab$, gdzie $a,b\in\mathbb F_2$.
Sumując cztery składniki, dostajemy
\[
 \sum_{a,b\in\mathbb F_2}
  (-1)^{a q(e_i)+bq(f_i)+ab}
     =2(-1)^{q(e_i)q(f_i)}.
\]
Iloczyn tych równości daje
$S(q)=2^r(-1)^{\sum_iq(e_i)q(f_i)}$.
Lewa strona nie używa bazy, więc prawa
wyznacza jednoznacznie Arfa. Dla sumy
ortogonalnej sumy $S$ mnożą się, dlatego
wykładniki znaków, czyli inwarianty Arfa,
dodają się modulo $2$.
\end{proof}

\begin{przyklad}[Dwie formy z tą samą macierzą przecięć]
\label{prz:Arf-dwie-22}
Niech $P=\Z e\oplus\Z f$ i
$\lambda(e,f)=1$, $\lambda(f,e)=-1$,
$\lambda(e,e)=\lambda(f,f)=0$.
Gdy $\mu(e)=\mu(f)=0$, para jest hiperboliczna
i $\operatorname{Arf}(q)=0$. Gdy
$\mu(e)=\mu(f)=1\in\Z/2$, otrzymujemy
$\operatorname{Arf}(q)=1$.
W drugim przypadku wszystkie trzy niezerowe
wektory płaszczyzny nad $\mathbb F_2$ mają
wartość $q=1$, bo
$q(e+f)=1+1+b(e,f)=1$.
Nie ma więc nawet jednowymiarowego
lagranżjanu z $q=0$. Macierz $\lambda$
jest w obu przykładach \emph{identyczna};
o przeszkodzie decyduje dodatkowa informacja
samoprzecięciowa.
\end{przyklad}

W tym wymiarze izomorfizm
$L_{4j+2}(\Z)\cong\Z/2$ przydziela formie
klasę $\operatorname{Arf}(q)$. Dowód, że
\emph{każda} forma z Arfem zerowym jest
stabilnie hiperboliczna, jest algebraiczną
częścią wspomnianego twierdzenia klasyfikacyjnego.

\section{Przykład z jedną nietrywialną pętlą}
\label{sec:L-pi1-Z-22}
\index{rozszczepienie Shanesona}

Pierwszym krokiem poza przypadek $\pi_1=1$
jest grupa nieskończenie cykliczna
$\pi_1=\langle t\rangle\cong\Z$. Jej pierścień
grupowy to pierścień wielomianów Laurenta
$\Z[\Z]=\Z[t,t^{-1}]$. Dla orientowalnej
rozmaitości charakter orientacji ma wartość
$w(t)=1$, a inwolucja używana w formach
przecięcia wysyła $t$ na $t^{-1}$.
Nie można tu zastąpić całego pierścienia
samymi współczynnikami całkowitymi:
wyrazy $t^k$ pamiętają drogi przecięcia
okrążające pętlę $k$ razy, jak w rachunku
$1-t$ w Sekcji~\ref{sec:pi-formy-22}.

\begin{twierdzenie}[Rozszczepienie Shanesona--Ranickiego;
wynik zewnętrzny]
\label{thm:Shaneson-splitting-22}
Niech $G$ będzie grupą z charakterem orientacji
$w:G\to\{\pm1\}$. Przedłużmy go na
$G\times\Z$ tak, aby generator nowego czynnika
zachowywał orientację. Dla $n\geq5$ grupy przeszkód
chirurgicznych spełniają
\begin{equation}\label{eq:Shaneson-splitting-22}
 L_n^s\bigl(\Z[G\times\Z],w\oplus1\bigr)
 \cong L_n^s\bigl(\Z[G],w\bigr)
       \oplus L_{n-1}^h\bigl(\Z[G],w\bigr).
\end{equation}
Drugi składnik ma dekorację $h$; dla dowolnego
$G$ nie należy jej bez dodatkowych hipotez
zamieniać na $s$.
\end{twierdzenie}

Jest to głęboki wynik o rozszczepianiu przeszkody
wzdłuż kierunku okręgu, a nie konsekwencja
formalnego rozpisania macierzy nad
$\Z[t,t^{-1}]$. Wariant geometryczny
\href{https://webhomes.maths.ed.ac.uk/~v1ranick/papers/shan.pdf}
{J.\ L.\ Shanesona, \emph{Wall's Surgery Obstruction
Groups for $G\times\Z$}, twierdzenie 5.1}
zakłada $n\geq7$, $G=\pi_1(K^{n-2})$ dla
zamkniętej gładkiej lub PL rozmaitości $K$
oraz $w=w_1(K)$. Ogólny izomorfizm grup z dekoracjami
$s$ i $h$ podaje
\href{https://webhomes.maths.ed.ac.uk/~v1ranick/papers/splitkl.pdf}
{A.\ Ranicki, \emph{Algebraic and Geometric Splittings
of the K- and L-Groups of Polynomial Extensions}, §6}.
Ranicki rozróżnia też geometryczne i algebraiczne
odwzorowania rozszczepienia; są one zgodne, gdy
$\operatorname{Wh}(G)=0$. Porównaj tabelę w
\href{https://webhomes.maths.ed.ac.uk/~v1ranick/books/scm.pdf}
{C.\ T.\ C.\ Wall, \emph{Surgery on Compact Manifolds},
twierdzenie 13A.8}.

Dla $G=1$ grupa Whiteheada jest zerowa
(Stwierdzenie~\ref{prop:wh-trywialna-20}),
więc dekoracje $s$ i $h$ po prawej stronie
\eqref{eq:Shaneson-splitting-22} się zgadzają.
Zatem z~\eqref{eq:L-Z-22} dostajemy
\begin{equation}\label{eq:L-ZZ-22}
 L_n^s\bigl(\Z[t,t^{-1}],1\bigr)
 \cong L_n(\Z)\oplus L_{n-1}(\Z).
\end{equation}
Rozpisanie dwóch składników, dla $n\geq5$,
daje konkretną tabelę:
\[
\begin{array}{c|cccc}
 n\bmod4&0&1&2&3\\ \hline
 L_n(\Z)&\Z&0&\Z/2&0\\
 L_{n-1}(\Z)&0&\Z&0&\Z/2\\ \hline
 L_n^s(\Z[t,t^{-1}],1)&\Z&\Z&\Z/2&\Z/2.
\end{array}
\]
W nieparzystych wymiarach, w których grupa
jednospójna znika, nowy składnik może więc
pozostać niezerowy. Na przykład dla $n=7$
rachunek daje
$L_7^s(\Z[t,t^{-1}])\cong
L_7(\Z)\oplus L_6(\Z)
=0\oplus\Z/2$;
jest to przesunięty składnik typu Arfa.
Dla $n=9$ otrzymujemy
$0\oplus L_8(\Z)\cong\Z$,
pochodzący z przesuniętego składnika sygnatury.

Nie jest to jeszcze obliczenie przeszkody
konkretnej mapy normalnej. Na przykład
$X=S^1\times S^6$ jest zamkniętą,
orientowalną rozmaitością wymiaru $7$
z $\pi_1(X)=\Z$, lecz jej mapa identycznościowa
z identycznościowymi danymi normalnymi ma
przeszkodę zerową: jest już prostą
równoważnością homotopijną.
Istnienie niezerowego elementu grupy $L_7^s$
nie oznacza, że przeszkoda tej mapy jest niezerowa.
Zbudowanie mapy normalnej o zadanej przeszkodzie
jest osobnym problemem realizacyjnym; w odpowiednim
zakresie wymiarów służą do tego twierdzenia Walla.
Podobnie tabela nie obejmuje generatora
odwracającego orientację, $w(t)=-1$;
algebraiczne grupy można badać także w małych
wymiarach, ale geometryczne kryterium Walla
stosujemy dopiero dla $n\geq5$.

\section{Przeszkoda normalnej mapy i twierdzenie Walla}
\label{sec:tw-Walla-22}

Rozważmy odwzorowanie normalne stopnia jeden
$(f,b):M^n\to X^n$ z
Sekcji~\ref{sec:jadra-chirurgii-21}.
Zakładamy, że $X$ jest zamkniętą gładką
rozmaitością lub skończonym kompleksem
Poincarégo z wybranymi danymi normalnymi.
Po chirurgiach \emph{poniżej} środka można
w dogodnych warunkach doprowadzić $f$ do
wysokiej spójności. W wymiarze parzystym
$n=2k$ pozostałe jądro nad $\Z[\pi]$
z formą $(\lambda_K,\mu_K)$ reprezentuje
\[
 \theta(f,b)=[K,\lambda_K,\mu_K]
       \in L^s_{2k}(\Z[\pi],w).
\]
Bez wcześniejszych chirurgii tę samą klasę
można zdefiniować na całym kompleksie łańcuchów
jako klasę kwadratowego kompleksu Poincarégo.
To ostatnie ujęcie obejmuje też wymiar nieparzysty.
Klasa nie zależy od wybranej kolejności
chirurgii: ślady tych zmian dają kobordyzm
odpowiednich kompleksów kwadratowych.

\begin{twierdzenie}[Twierdzenie Walla o przeszkodzie chirurgicznej]
\label{thm:Wall-L-22}
\index{twierdzenie Walla@Twierdzenie Walla}
Niech $M^n$ i $X^n$ będą spójnymi, zamkniętymi,
gładkimi rozmaitościami, $n\geq5$, i niech
$(f,b):M\to X$ będzie odwzorowaniem normalnym
stopnia jeden z Definicji~\ref{def:normal-map-21},
z danymi normalnymi nad pewną stabilną wiązką
$\xi\to X$ jak w~\eqref{eq:normal-map-21}.
Oznaczmy $\pi=\pi_1(X)$ i przez $w$ charakter
orientacji. Wtedy przeszkoda
$\theta(f,b)\in L^s_n(\Z[\pi],w)$ jest niezmiennikiem
normalnej klasy kobordyzmu.
Równość $\theta(f,b)=0$ zachodzi wtedy i tylko wtedy,
gdy $(f,b)$ jest normalnie kobordantne z prostą
równoważnością homotopijną $f':M'\to X$.
W wersji zwartej z brzegiem zakładamy ponadto,
że odwzorowanie na brzegu jest prostą równoważnością
homotopijną i wykonujemy kobordyzm oraz chirurgie
relatywnie do brzegu.
\end{twierdzenie}

\begin{twierdzenie}[Wersja $h$ i porównanie dekoracji;
 wyniki zewnętrzne]\label{thm:Wall-h-22}
Przy tych samych założeniach o wymiarze $n\geq5$
i danych normalnych, ale z wymaganiem jedynie
\emph{zwykłej} równoważności homotopijnej na brzegu,
istnieje przeszkoda
\[
 \theta^h(f,b)\in L^h_n(\Z[\pi],w).
\]
Zależy ona tylko od normalnej klasy kobordyzmu
relatywnie do brzegu i znika dokładnie wtedy,
gdy tę mapę można przez taki kobordyzm zastąpić
równoważnością homotopijną. Dla pustego brzegu
warunek brzegowy jest oczywiście zbędny.

Jeżeli mapa brzegowa jest prosta, zapomnienie
warunku prostoty daje homomorfizm
$j:L^s_n(\Z[\pi],w)\to L^h_n(\Z[\pi],w)$ oraz
\[
 \theta^h(f,b)=j\bigl(\theta^s(f,b)\bigr).
\]
Jądro i koniec obrazu $j$ opisuje fragment
\emph{ciągu dokładnego Rothenberga}
\begin{equation}\label{eq:Rothenberg-sh-22}
 \cdots\longrightarrow
 \widehat H^{\,n+1}\bigl(\Z/2;\operatorname{Wh}(\pi)\bigr)
 \longrightarrow L^s_n(\Z[\pi],w)
 \xrightarrow{\ j\ } L^h_n(\Z[\pi],w)
 \longrightarrow
 \widehat H^{\,n}\bigl(\Z/2;\operatorname{Wh}(\pi)\bigr)
 \longrightarrow\cdots .
\end{equation}
W grupach Tate'a działa inwolucja dualności
wyznaczona przez $w$ na grupie Whiteheada.
\end{twierdzenie}

Twierdzenie o przeszkodzie w dekoracji $h$
oraz ciąg~\eqref{eq:Rothenberg-sh-22} są tutaj
przyjętymi wynikami zewnętrznymi:
\href{https://www.maths.ed.ac.uk/~v1ranick/surgery/ranicki.pdf}
{Ranicki, \emph{An Introduction to Algebraic Surgery},
 wstęp i przykład 2.16} podaje kryterium
znikania przeszkody oraz warunek brzegowy;
\href{https://www.maths.ed.ac.uk/~v1ranick/books/scm.pdf}
{Wall, \emph{Surgery on Compact Manifolds}, §17D},
odsyła w sprawie ciągu do
\href{https://webhomes.maths.ed.ac.uk/~v1ranick/papers/shan.pdf}
{Shanesona, \emph{Wall's Surgery Obstruction
Groups for $G\times\Z$}, propozycja 4.1}.
Wall sam nazywa swoje porównanie w §17D
zarysem, więc nie używamy go jako pełnego dowodu.
Dokładność~\eqref{eq:Rothenberg-sh-22} pokazuje,
że zniknięcie przeszkody $h$ nie musi usuwać
warunku prostej torsji potrzebnego do wersji $s$.
Gdy $\operatorname{Wh}(\pi)=0$, grupy Tate'a
znikają i $j$ jest izomorfizmem; w szczególności
dotyczy to $\pi=1$.

\begin{przyklad}[Grupy Tate'a dla $C_5$]
\label{prz:Tate-C5-22}
\index{grupa Tate'a!przykład dla $C_5$}
Wróćmy do $R=\Z[C_5]$ i jednostki
$u=1-t-t^{-1}$ z Przykładu~\ref{prz:dekoracje-C5-22}.
\href{https://webhomes.maths.ed.ac.uk/~v1ranick/papers/milnorwh.pdf}
{J.\ Milnor, \emph{Whitehead Torsion}, przykład 6.6}
podaje $\operatorname{Wh}(C_5)\cong\Z$ oraz
generator $[t+t^{-1}-1]=[-u]=[u]$;
ostatnia równość zachodzi, bo $[-1]=0$
w grupie Whiteheada. Ponieważ $\bar u=u$,
inwolucja dualności działa na generatorze
identycznościowo, a więc działa tak na całej
grupie $\operatorname{Wh}(C_5)$.

Oznaczmy tę inwolucję przez $T$. Dla
$A=\operatorname{Wh}(C_5)=\Z\cdot[u]$
mamy $1-T=0$ oraz $1+T=2$. Ze wzorów
definiujących kohomologię Tate'a otrzymujemy
\[
 \widehat H^{2j}(\Z/2;A)
 =\frac{\ker(1-T)}{\operatorname{im}(1+T)}
 =\frac{A}{2A}\cong\Z/2,
 \qquad
 \widehat H^{2j+1}(\Z/2;A)
 =\frac{\ker(1+T)}{\operatorname{im}(1-T)}=0.
\]
Niezerową klasą parzystą jest reszta $[u]$
modulo $2A$. Po podstawieniu do
ciągu~\eqref{eq:Rothenberg-sh-22}, dla
$n\geq6$ parzystego, dostajemy fragment
\[
 0\longrightarrow L^s_n(R,1)
 \xrightarrow{j}L^h_n(R,1)
 \longrightarrow\Z/2\longrightarrow L^s_{n-1}(R,1),
\]
a dla $n\geq5$ nieparzystego
\[
 \Z/2\longrightarrow L^s_n(R,1)
 \xrightarrow{j}L^h_n(R,1)\longrightarrow0.
\]
W parzystym wymiarze $j$ jest zatem injekcją,
a w nieparzystym surjekcją. Dokładność przy
parzystym $n\geq6$ daje ponadto
\[
 0\longrightarrow\operatorname{coker}j_n
 \longrightarrow\Z/2\longrightarrow
 \ker j_{n-1}\longrightarrow0.
\]
Zatem przynajmniej jedna z dwóch sąsiednich map
$j_n,j_{n-1}$ nie jest izomorfizmem, ale rachunek
nie wyznacza, która. Nie rozstrzyga też grup
$L^s_n,L^h_n$ jako grup abstrakcyjnych ani
przeszkody konkretnej mapy normalnej.
\end{przyklad}

\noindent\textbf{Zakres wyprowadzenia.}
Najpierw rozpiszemy argument dostateczności
dla $\pi=1$ i $n=2k\geq6$, a potem
przeprowadzimy wersję nad $\Z[\pi]$.
W wymiarze nieparzystym $n=2k+1\geq5$
zbudujemy przeszkodę z dwóch lagranżjanów
i pokażemy, jak jej zanik prowadzi do chirurgii.
Potrzebne wejścia geometryczne
pochodzą z
\href{https://webhomes.maths.ed.ac.uk/~v1ranick/books/scm.pdf}
{C.~T.~C. Walla, \emph{Surgery on Compact Manifolds},
§§1, 5--6, szczególnie twierdzenia 1.2, 5.2,
5.6, 6.4 oraz lemat 5.7}.
Lemat 5.3 o prostym lagranżjanie dowiedziemy
poniżej; inne użyte wyniki Walla oznaczamy
wprost z ich hipotezami. Wersję z dekoracją
$h$ i zwykłą równoważnością homotopijną
traktujemy oddzielnie jako wynik zewnętrzny;
§17D Walla podaje jedynie szkic porównania
tej dekoracji z wersją $s$.

\noindent\textbf{Cztery kroki dowodu.}
Są to cztery odrębne kroki, a nie samo
wydłużenie rachunku form nad $\Z$:
\begin{enumerate}[label=(\arabic*)]
\item \emph{Chirurgie poniżej środka.}
  Poniżej konstruujemy zgodne obramowanie
  dla pojedynczej klasy w zakresie $2r<n$.
  Skończona indukcja po względnych komórkach
  i kontrola nowych jąder jest treścią
  twierdzenia 1.2 Walla.
\item \emph{Realizacja w środku.}
  W wymiarze parzystym wyliczymy etykiety przecięć i zbudujemy
  bazę hiperboliczną. Zanurzenie sfery w środku
  z warunku $\mu=0$ wymaga obramowanych dysków
  Whitneya; dla wielu sfer trzeba usuwać
  również ich wzajemne przecięcia zgodnie
  z formą $\lambda$ i udoskonaleniem $\mu$;
  używamy tu twierdzenia 5.2 Walla.
\item \emph{Prosta równoważność i niezmienniczość.}
  Zmiana kolejności chirurgii musi dawać
  tę samą klasę przeszkody. Dla ogólnej
  grupy podstawowej wymaga to bazowanych
  kompleksów nad $\Z[\pi]$, kontroli torsji
  Whiteheada i analizy kobordyzmu normalnego.
  Poniżej wyprowadzimy z ciągu dokładnego
  warunek lagranżjanu i rachunek znoszenia
  przecięć na brzegu kobordyzmu. Przygotowanie
  kobordyzmu wymaga twierdzenia 1.4 Walla;
  bazowość, prosta dualność i realizacja
  względnych immersji są osobnymi wejściami
  z jego §§1--2.
\item \emph{Wymiar nieparzysty.}
  Jądra dwóch sąsiednich wymiarów są mierzone
  przecięciem i ilorazem dwóch lagranżjanów
  formacji. Poniżej wyprowadzimy jej ciąg
  dokładny; relacja stabilna i realizacja
  odpowiadających jej ruchów uchwytów
  podamy z jawnym użyciem §6 Walla.
\end{enumerate}
Rachunek poniżej dowodzi kryterium po przyjęciu
wskazanych lematów geometrycznych Walla;
ich własnych dowodów nie przypisujemy temu tekstowi.

\begin{lemat}[Algebraiczny efekt usunięcia jednej pary]
\label{lem:para-Wall-22}
Niech $(K,\lambda,\mu)$ będzie sumą $r$ płaszczyzn
hiperbolicznych nad $\Z$, z bazą
$e_1,f_1,\ldots,e_r,f_r$ taką, że
$\lambda(e_i,f_j)=\delta_{ij}$,
$\lambda(e_i,e_j)=\lambda(f_i,f_j)=0$
oraz $\mu(e_i)=\mu(f_i)=0$.
Wówczas
\[
 e_r^\perp/\Z e_r
 \cong\langle e_1,f_1,\ldots,e_{r-1},f_{r-1}\rangle
\]
jako formy kwadratowe.
\end{lemat}
\begin{proof}
Zapiszmy
$x=\sum_i(a_i e_i+b_i f_i)$.
Z definicji bazy hiperbolicznej
$\lambda(e_r,x)=b_r$. Stąd $x\in e_r^\perp$
dokładnie wtedy, gdy $b_r=0$.
Pozostały współczynnik $a_r$ znika po
przejściu do ilorazu przez $\Z e_r$.
Pary o indeksach $i<r$ są ortogonalne
do $e_r,f_r$, więc ich $\lambda$ i $\mu$
przechodzą do ilorazu bez zmiany.
\end{proof}

\begin{lemat}[Od prostego lagranżjanu do bazy hiperbolicznej]
\label{lem:lagr-hip-Z-22}
Niech $(P,\lambda,\mu)$ będzie nieosobliwą
formą $\epsilon$-kwadratową nad
$\Lambda=\Z[\pi]$, $\epsilon=(-1)^k$.
Załóżmy, że $P$ jest wolny z preferowaną bazą,
a $L\subset P$ jest wolnym lagranżjanem z bazą
$(e_1,\ldots,e_r)$, która przedłuża się do
preferowanej bazy $P$. Niech ponadto
$P/L\to L^*$, $[x]\mapsto\lambda(-,x)|_L$,
będzie prostym izomorfizmem.
Wtedy istnieje preferowana baza
$(e_1,f_1,\ldots,e_r,f_r)$, w której
$\lambda(e_i,f_j)=\delta_{ij}$ oraz
$\lambda(e_i,e_j)=\lambda(f_i,f_j)=0$ i
$\mu(e_i)=\mu(f_i)=0$.
\end{lemat}
\begin{proof}
Wybierzmy podniesienia $f_j^0\in P$ dualnej
bazy $L^*$, zatem
$\lambda(e_i,f_j^0)=\delta_{ij}$.
Jądrem $P\to L^*$ jest $L^\perp=L$,
więc klasy $f_j^0$ są bazą $P/L$.
Założenie o prostocie zapewnia, że
$(e_i,f_j^0)$ reprezentuje preferowaną
klasę baz $P$.

Połóżmy $B_{ij}=\lambda(f_i^0,f_j^0)$
i wybierzmy $m_i\in\Lambda$ reprezentujące
$\mu(f_i^0)$. Zdefiniujmy
\begin{equation}\label{eq:poprawka-lagr-Wall-22}
 f_i=f_i^0-m_i e_i-
       \sum_{j>i}B_{ij}e_j.
\end{equation}
Dodajemy wyłącznie wektory z $L$, toteż
$\lambda(e_i,f_j)=\delta_{ij}$.
Dla $i<j$ składnik $-B_{ij}e_j$ w $f_i$
znosi pierwotne parowanie:
\[
 \lambda(f_i,f_j)=B_{ij}-B_{ij}=0.
\]
Przy $i>j$ używamy hermitowskości.
Ponieważ $\lambda(f_i^0,e_j)=0$ dla $j>i$
i $\lambda(L,L)=0$, równanie
\eqref{eq:quad-add-22} daje
\[
 \mu(f_i)=[m_i]+[\lambda(f_i^0,-m_i e_i)]
   =[m_i-\epsilon\bar m_i]=0.
\]
Równanie~\eqref{eq:quad-diag-22} daje teraz
$\lambda(f_i,f_i)=0$.
Zmiana $f_i^0\mapsto f_i$ jest elementarną
zmianą bazy i nie zmienia jej prostej klasy.
To dowodzi tezy również przy nieprzemiennym
$\Lambda$; dla $\Lambda=\Z$ odzyskujemy
poprzedni rachunek parzystych macierzy
i poprawki Arfa.
\end{proof}

\begin{stwierdzenie}[Przeciwna forma daje odwrotność]
\label{prop:odwrotna-forma-Wall-22}
Jeśli $(P,\lambda,\mu)$ jest prostą
nieosobliwą formą $\epsilon$-kwadratową
na wolnym bazowanym module, to
\[
 (P,\lambda,\mu)\oplus(P,-\lambda,-\mu)
\]
jest formą hiperboliczną. Stąd jej druga
składowa reprezentuje klasę przeciwną
w $L^s_{2k}(\Lambda,w)$.
\end{stwierdzenie}
\begin{proof}
Weźmy diagonalny podmoduł
$\Delta=\{(x,x):x\in P\}\subset P\oplus P$.
Na $\Delta$ obie formy znoszą się:
$\lambda((x,x),(y,y))=0$ oraz
$\mu((x,x))=0$. Jest to bezpośredni
składnik z bazą $(e_i,e_i)$, jeśli
$(e_i)$ jest bazą $P$. Iloraz
$(P\oplus P)/\Delta$ identyfikujemy
z $P$ przez $[(x,y)]\mapsto x-y$.
Parowanie z $\Delta$ jest, z dokładnością
do znaku, izomorfizmem dualności
$P\to P^*$ wyznaczonym przez $\lambda$.
Jest proste z założenia. Lemat
\ref{lem:lagr-hip-Z-22} pokazuje więc,
że $\Delta$ daje bazę hiperboliczną.
\end{proof}

\begin{lemat}[Dlaczego obramowanie istnieje poniżej środka]
\label{lem:ramy-ponizej-srodka-Wall-22}
Niech $(f,b):M^{2k}\to X^{2k}$ będzie
odwzorowaniem normalnym, $k\geq3$, a
$\alpha\in\pi_{r+1}(f)$, gdzie $1\leq r\leq k-1$.
Klasa $\alpha$ ma reprezentanta będącego
zanurzoną $r$-sferą z obramowaniem zgodnym
z danymi normalnymi. Chirurgia na tej sferze
zabija $\alpha$ i nie przywraca jąder
w stopniach mniejszych od $r+1$.
\end{lemat}
\begin{proof}
Reprezentant klasy względnej jest parą
$u:S^r\to M$ i $F:D^{r+1}\to X$ z
$F|_{S^r}=f\circ u$. Ponieważ
$2r<2k$, ogólne położenie pozwala
zastąpić $u$ homotopicznym zanurzeniem:
oczekiwany wymiar jego samoprzecięć
wynosi $2r-2k<0$.

Połóżmy $q=2k-r$. Stabilne dane normalne
$b$ utożsamiają $u^*\nu_M$ z
$(f\circ u)^*\xi$ po stabilizacji.
Ta ostatnia wiązka przedłuża się przez
$F$ nad dyskiem, więc jest stabilnie
trywialna. Ponieważ
$u^*TM\oplus u^*\nu_M$ jest stabilnie
trywialna oraz $TS^r\oplus\R$
jest trywialna, wzór
$u^*TM=TS^r\oplus\nu_u$ pokazuje,
że również $\nu_u$ ma zgodne stabilne
obramowanie. Ponieważ
$q\geq r+2$, można usunąć stabilizację.
Istotnie, $O(q+1)/O(q)\cong S^q$;
długi ciąg homotopii i $(q-1)$-spójność
$S^q$ pokazują, że $O(q)\to O$
jest izomorfizmem na $\pi_j$ dla
$j\leq r$. Dotyczy to zarówno przeszkody
do trywialności $\nu_u$ nad $S^r$,
jak i różnicy dwóch jej obramowań.
Wybieramy rzeczywiste obramowanie.
Jeżeli po zestawieniu go z różniczką
$du:TS^r\to u^*TM$ jego stabilna klasa
różni się od ramy wyznaczonej przez $b$
i $F$ o element $\pi_r(O)$, podnosimy
ten element do $\pi_r(O(q))$ i skręcamy
obramowanie o jego odwrotność.
Otrzymujemy dokładnie zadaną stabilną
klasę ramy, a nie tylko dowolną
trywializację $\nu_u$.

Dysk $F$ i zgodne obramowanie przedłużają
normalne dane przez ślad, zgodnie z
Twierdzeniem~\ref{thm:przedluzenie-mapy-21}.
Ślad ma względem $M$ jeden uchwyt indeksu
$r+1$, który zabija $\alpha$.
Względem nowego brzegu indeks wynosi
$2k-r\geq k+1>r+1$; dlatego ciąg
homotopii pary nie wprowadza nowych
jąder poniżej stopnia $r+1$.
\end{proof}

Przypadek $r=0$ odpowiada dołączeniu uchwytu
indeksu jeden wzdłuż dwóch punktów. Jeżeli
$M$ ma kilka składowych, takimi uchwytami
możemy najpierw je połączyć, przedłużając
mapę wzdłuż wybranych dróg w spójnym $X$.
Ramy obu dysków zaczepienia dobieramy
zgodnie z danymi normalnymi $b$ dokładnie
tak jak w przypadku $h_i=1$ poniżej.

\begin{stwierdzenie}[Skończona redukcja poniżej środka;
 szczególny przypadek twierdzenia 1.2 Walla]
\label{prop:skonczona-redukcja-Wall-22}
Niech $(f,b):M^{2k}\to X^{2k}$ będzie normalnym
odwzorowaniem zwartej rozmaitości gładkiej
bez brzegu,
$k\geq3$, a $X$ niech będzie spójną
przestrzenią o typie skończonego kompleksu
CW z danymi normalnymi wymaganymi dla $(f,b)$.
Istnieje skończony
ciąg chirurgii na sferach wymiarów mniejszych
od $k$, po którym otrzymujemy normalne
odwzorowanie $f':M'\to X$ będące
$k$-spójne. Wszystkie użyte uchwyty mają
indeks co najwyżej $k$.
\end{stwierdzenie}
\begin{proof}
Po ewentualnym połączeniu składowych zakładamy,
że $M$ jest spójna. Z
Twierdzenia~\ref{thm:cw-morse-17} wynika,
że $M$ ma skończony model CW: wybierzmy
równoważność $u:K\to M$ ze skończonym $K$.
Wybieramy również skończony model $L\simeq X$.
Przenosimy $f\circ u$ na mapę $K\to L$,
przybliżamy ją komórkowo i tworzymy
podwójny walec odwzorowań
$M\leftarrow K\to L$. Ma on skończoną
względną strukturę CW nad $M$ i typ
homotopijny zwykłego walca $f$;
oznaczmy ten model przez $C_f$. Wybieramy
jego strukturę komórkową tak, by
$M\hookrightarrow C_f$ było inkluzją,
$C_f\simeq X$, a jego względna struktura
komórkowa nie miała komórek wymiaru zero.
Względne wierzchołki łączymy z $M$
lasem krawędzi, którego każda składowa
spotyka $M$ w jednym wierzchołku,
i zwijamy ten las względem $M$. Wybieramy
równoważność $C_f\to X$, której złożenie
z inkluzją $M$ jest homotopijne z $f$;
dane normalne przenosimy wzdłuż tej homotopii.
Porządkujemy względne komórki wymiarów
co najwyżej $k$ według wymiaru i piszemy
\[
 C_0=M,\qquad C_i=C_{i-1}\cup e^{h_i},
 \qquad 1\leq h_i\leq k.
\]

Indukcyjnie budujemy ślad $W_i$ chirurgii
od $M$ do $M_i$ i równoważność homotopijną
$W_i\simeq C_i$ zgodną z odwzorowaniem do
$C_f$. Dla $i=0$ bierzemy $W_0=M\times I$.
Załóżmy, że zbudowaliśmy $W_{i-1}$.
Komórka $e^{h_i}$ ma odwzorowanie dołączające
$S^{h_i-1}\to C_{i-1}$ i klasę
charakterystyczną w
$\pi_{h_i}(C_i,C_{i-1})$. Dotychczasowe
uchwyty mają indeks co najwyżej $k$.
Po odwróceniu śladu względem $M_{i-1}$
ich indeksy wynoszą co najmniej
\[
 (2k+1)-k=k+1>h_i.
\]
Stąd para $(W_{i-1},M_{i-1})$ jest
$k$-spójna. Po zastąpieniu rozważanych
odwzorowań walcami odwzorowań i użyciu
$W_{i-1}\simeq C_{i-1}$ otrzymujemy,
ponieważ $h_i\leq k$, izomorfizm względnych
klas homotopii
\[
 \pi_{h_i}(C_f,M_{i-1})
 \xrightarrow{\ \cong\ }
 \pi_{h_i}(C_f,C_{i-1}).
\]
Podnosimy nim klasę dysku charakterystycznego
$e^{h_i}$, a nie tylko klasę jego brzegu.
Reprezentant podniesionej klasy składa się
ze sfery $u_i:S^{h_i-1}\to M_{i-1}$
i dysku $D^{h_i}\to C_f$ ją wypełniającego.
Po złożeniu z $C_f\simeq X$ otrzymujemy
klasę względną odwzorowania $f_{i-1}$.
Zgodność podniesienia z klasą
charakterystyczną mówi ponadto, że obraz
$u_i$ w $C_{i-1}$ jest homotopijny
z odwzorowaniem dołączającym komórkę.

Dla $h_i\geq2$
Lemat~\ref{lem:ramy-ponizej-srodka-Wall-22}
realizuje tę klasę przez zanurzoną
obramowaną sferę $S^{h_i-1}$ i chirurgię
z przedłużeniem danych normalnych.
Dla $h_i=1$ sfera $S^0$ składa się
z dwóch punktów, a podniesiony dysk
jest drogą łączącą ich obrazy w $C_f$.
Wiązka normalna nad tą drogą jest
trywialna. Dobieramy małe dyski wokół
obu punktów z ramami zgodnymi ze stabilnymi
danymi $b$; jest to możliwe, bo
$\pi_0 O(2k)\to\pi_0 O$ jest izomorfizmem.
Dołączamy wzdłuż nich uchwyt indeksu jeden
i przedłużamy dane normalne, jak w przypadku
$r=0$ twierdzenia 1.1 Walla.

Nowy ślad powstaje z $W_{i-1}$ przez
dołączenie uchwytu indeksu $h_i$.
Zgodność klasy dołączającej z klasą
komórki i homotopijna niezmienniczość
dołączania komórek dają
$W_i\simeq C_i$ nad $C_f$.

Po ostatniej takiej komórce otrzymujemy
$C_j$. Para $(C_f,C_j)$ nie ma komórek
względnych wymiaru co najwyżej $k$,
więc jest $k$-spójna. Para $(W_j,M_j)$
też jest $k$-spójna, bo wszystkie jej
odwrócone uchwyty mają indeks co najmniej
$k+1$. Zatem złożenie
$M_j\to W_j\simeq C_j\hookrightarrow C_f$
jest $k$-spójne: pierwszy i ostatni człon
mają tę własność przez odpowiednie
ciągi homotopii par, a środkowy jest
równoważnością. Równoważność $C_f\simeq X$
pokazuje, że $f_j:M_j\to X$ jest
$k$-spójne.
Skończoność ciągu wynika ze skończoności
wybranego modelu CW.
\end{proof}

Dowód redukcji używa pojedynczego kroku
geometrycznego z Lematu~\ref{lem:ramy-ponizej-srodka-Wall-22}
i schematu indukcji z
\href{https://webhomes.maths.ed.ac.uk/~v1ranick/books/scm.pdf}
{Walla, §1, twierdzenie 1.2}. Ogólna wersja
obejmuje także inne wymiary i kategorię PL.

\begin{proof}[Rozwinięcie dowodu twierdzenia Walla w przypadku
 $\pi=1$, $n=2k\geq6$]
\textbf{Krok 1: chirurgie poniżej środka.}
Normalne dane $b$ pozwalają przedłużać
odwzorowanie nad śladami chirurgii, jak w
Twierdzeniu~\ref{thm:przedluzenie-mapy-21}.
Lemat~\ref{lem:ramy-ponizej-srodka-Wall-22}
konstruuje zgodne zanurzenie i obramowanie
dla jednej klasy, a
Stwierdzenie~\ref{prop:skonczona-redukcja-Wall-22}
uzasadnia skończone powtarzanie tego kroku.
Po tych operacjach jedynym możliwym
jądrem chirurgii jest środkowe $K_k$.
Identyfikacja $K_k(M)\cong\pi_{k+1}(f)$
i prosta dualność Poincarégo dają
na nim nieosobliwą formę
$(K_k,\lambda,\mu)$ z
Sekcji~\ref{sec:kwadratowa-22}.
Wersja nad $\Lambda$ wymaga dodatkowo
preferowanej stabilnej bazy jądra:
to dokładnie lemat 5.1 Walla.
Podkreślamy, że warunek kwadratowy dotyczy
klas z \emph{normalnym} obramowaniem;
dla dowolnej sfery bez tych danych wzór
na samoprzecięcie ma dodatkowy składnik
Eulera wiązki normalnej (Wall, §5.2).

\textbf{Krok 2: przejście od klasy zerowej do
bazy hiperbolicznej.}
Równość $[K]=0$ w $L_{2k}(\Z)$
oznacza, że dla pewnych $r,s$
forma $K\oplus H_\epsilon^r$ jest
izomorficzna z $H_\epsilon^s$.
Równoważnie, po stabilizacji ma
lagranżjan; przejście od niego do
rzeczywistej bazy z parami $e_i,f_i$
rozpisuje Lemat~\ref{lem:lagr-hip-Z-22}.
Każdą dodawaną
płaszczyznę realizujemy \emph{lokalną}
chirurgią na trywialnej obramowanej
$(k-1)$-sferze: nowa rozmaitość jest
$M\#(S^k\times S^k)$, a dwie sfery
czynnikowe przecinają się raz z
odpowiednim znakiem i mają $\mu=0$.
Ślad dostarcza normalnego kobordyzmu.
To jest lemat 5.5 Walla; wyjaśnia, dlaczego
stabilizacja algebraiczna jest dozwolonym
ruchem geometrycznym. Otrzymujemy bazę
$e_i,f_i$ jak w Lemacie~\ref{lem:para-Wall-22}.

\textbf{Krok 3: jedna rzeczywista chirurgia
w środku.}
Warunek $\mu(e_r)=0$ i
Stwierdzenie~\ref{prop:osadzenie-Wall-52-22},
używające obramowanego triku Whitneya
przy $k\geq3$, dają zanurzoną obramowaną
sferę $S^k\subset M$ reprezentującą $e_r$
oraz dysk w $X$ wypełniający jej obraz.
Wykonujemy chirurgię nad $X$.
Jej ślad $W$ powstaje z $M\times I$
przez dołączenie uchwytu indeksu $k+1$.
Oznaczmy nowy brzeg przez $M'$ i piszmy
$K_i(W)$ dla jądra mapy śladu do $X\times I$.
Oba względne kompleksy uchwytowe mają po
jednym generatorze:
\[
 H_{k+1}(W,M;\Z)\cong\Z,\qquad
 H_k(W,M';\Z)\cong\Z.
\]
W pierwszym ciągu dokładnym mapa brzegu
$\Z\to K_k(M)$ wysyła $1$ na $e_r$.
Jest injektywna, gdyż partner $f_r$ daje
$\lambda(e_r,f_r)=1$, więc $e_r$ jest
prymitywny. Otrzymujemy
\[
 0\longrightarrow\Z e_r\longrightarrow K_k(M)
 \longrightarrow K_k(W)\longrightarrow0.
\]
Od strony $M'$ uchwyt ma indeks $k$.
Złożenie $K_k(M)\to K_k(W)\to
H_k(W,M';\Z)\cong\Z$ mierzy przecięcie
z $e_r$. Dla lewych modułów nad
$\Lambda$ odpowiada mu \emph{liniowy}
homomorfizm
\[
 \ell_{e_r}:K_k(M)\longrightarrow\Lambda,
 \qquad \ell_{e_r}(z)=\epsilon\lambda(z,e_r),
 \quad \epsilon=(-1)^k.
\]
Istotnie, $\ell_{e_r}(az)=a\ell_{e_r}(z)$,
a hermitowskość pokazuje, że jego jądrem
jest $e_r^\perp$. Przy $\Lambda=\Z$ jest to
ten sam wzór co $\lambda(e_r,z)$.
Ponieważ $\lambda(e_r,e_r)=0$, mamy
$\ell_{e_r}(e_r)=0$; mapa przechodzi więc
przez $K_k(W)=K_k(M)/\Z e_r$.
Drugi ciąg dokładny daje
\[
 0\longrightarrow K_k(M')\longrightarrow K_k(W)
 \xrightarrow{\,\ell_{e_r}\,}\Z
 \longrightarrow K_{k-1}(M')\longrightarrow0.
\]
Ponieważ $f_r$ przechodzi na $1$,
$K_{k-1}(M')=0$, a
$K_k(M')\cong e_r^\perp/\Z e_r$.
Forma i udoskonalenie kwadratowe przechodzą
na ten iloraz, gdyż $\mu(e_r)=0$.
W ogólnym przypadku wszystkie wystąpienia
$\Z$ w tych ciągach zastępuje $\Lambda$.
Bez partnera o przecięciu $1$ mógłby
powstać niezerowy $K_{k-1}(M')$.
Lemat~\ref{lem:para-Wall-22} usuwa dokładnie
parę $e_r,f_r$. Powtarzamy ten krok $r$ razy.

\textbf{Krok 4: otrzymanie równoważności.}
Po usunięciu wszystkich par $K_k=0$,
a inne jądra zostały usunięte w Kroku 1
lub znikają przez dualność. Zatem nowa
mapa stopnia jeden jest izomorfizmem
homologii. Ponieważ obie strony są
jednospójne i mają typ kompleksów CW,
twierdzenie Whiteheada
(Twierdzenie~\ref{thm:whitehead})
daje równoważność homotopijną.
Dla $\pi=1$ grupa Whiteheada jest zerowa,
więc jest to równoważność prosta.

\textbf{Konieczność i niezależność wyborów.}
Przeszkoda prostej równoważności wynosi
zero, bo jej kompleks jądra jest ściągalny.
Rozpiszmy rachunek z lematu~5.7 Walla,
który daje niezależność od normalnego
kobordyzmu. Niech ogólniej
$\Phi:(N^{2k+1},M^{2k})\to(Y,X)$ będzie
odwzorowaniem normalnym stopnia jeden
do spójnej prostej pary Poincarégo; wolno
tu dopuścić niespójny brzeg $M$.
W całym rachunku współczynnikiem jest
$\Lambda=\Z[\pi_1(Y)]$, z inwolucją
orientacyjną celu $Y$. Na ewentualnie
niespójnym $M$ formy
$\lambda$ i $\mu$ liczymy na nakryciu
indukowanym z uniwersalnego nakrycia $Y$,
jak w uwadze~5.2.1 Walla. Zakładamy,
że $M\to X$ i $N\to Y$ są $k$-spójne,
$K_k(N,M)=0$, a preferowana stabilna baza
$K_{k+1}(N,M)$ jest reprezentowana przez
rzeczywistą bazę. To dokładne hipotezy
lematu~5.7; samo ich osiągnięcie jest
oddzielnym zadaniem chirurgicznym.

\emph{Krok pierwszy: ciąg i jego bazy.}
Połóżmy
\[
 A=K_{k+1}(N,M),\qquad P=K_k(M),
 \qquad C=K_k(N).
\]
Spójność usuwa niższe jądra, a względna
dualność pozostałe wyższe. Przy
$K_k(N,M)=0$ ciąg długi pary redukuje się do
\begin{equation}\label{eq:wall-57-ciag-22}
 0\longrightarrow A\xrightarrow{\,\partial\,}P
 \xrightarrow{\,i_*\,}C\longrightarrow0.
\end{equation}
Jest to \emph{bazowany} ciąg prostych modułów:
to osobna treść wniosku z lematu~2.4 Walla,
nie własność dowolnego krótkiego ciągu
dokładnego. Wersja dualności Poincarégo
dla kompleksów jądra, lemat~2.6 Walla,
daje prosty izomorfizm $C\cong A^*$.
W lemacie~2.4 bierze się bazowany krótki ciąg
\emph{skończonych kompleksów wolnych} otrzymany
z trzech kompleksów jądra map $M\to X$,
$N\to Y$ i $(N,M)\to(Y,X)$. W naszej sytuacji
wniosek z lematu~2.4 wymaga zaniku
$K_i(N,M)$ poza stopniem $k+1$ i $K_i(N)$
poza stopniem $k$. Ponadto po zmianie
współczynników na dowolny moduł $B$ nad
$\Lambda$ muszą znikać
$K_{k+2}(N,M;B)$ i $K_{k+1}(N;B)$;
oznaczamy tak homologie tych samych
kompleksów jądra ze współczynnikami $B$.
Te hipotezy
zapewniają, że jedyne pozostałe moduły homologii
są stabilnie wolne i mają preferowane klasy baz.
Addytywność torsji dla tego ciągu daje zerową
torsję ciągu~\eqref{eq:wall-57-ciag-22}.
Lemat~2.6 stosuje \emph{prostą} dualność pary
Poincarégo do tych samych kompleksów. Po
redukcji kompleksu jądra $N\to Y$ do jednego
stopnia jego kohomologia to $C^*$; iloczyn
z klasą fundamentalną identyfikuje ją prosto
z $A$. Ponieważ $A$ i $C$ są stabilnie wolne,
dualizacja daje $C\cong A^*$.
Sam
izomorfizm zwykłych grup homologii nie dawałby
ani preferowanych baz, ani prostoty dualności.
Jego wzór w naszej konwencji lewych
modułów rozpoznamy po Kroku 2. Na razie
na całym $P$ określmy
\[
 \widetilde\rho:P\longrightarrow A^*,
 \qquad \widetilde\rho(z)(x)=\lambda(\partial x,z).
\]
Na $A^*=\operatorname{Hom}_\Lambda(A,\Lambda)$
działanie lewe określa
$(a\cdot h)(x)=h(x)\bar a$.
Dlatego $\widetilde\rho(az)=a\cdot
\widetilde\rho(z)$;
strzałka z $\lambda(z,\partial x)$
byłaby przy tej konwencji antyliniowa.

\emph{Krok drugi: znikanie formy.}
Weźmy $x,y\in A$. Obramowane immersje
reprezentujące $\partial x,\partial y$
na $M$ można przedłużyć do obramowanych
immersji $T_x,T_y$ w $N$, których dziedziną
jest suma skończona $(k+1)$-sfer z usuniętymi
dyskami. Zniknięcie klasy brzegowej w $C$
wynika z~\eqref{eq:wall-57-ciag-22};
przejście od tego faktu homologicznego do
takiego \emph{obramowanego} przedłużenia
korzysta ze względnego twierdzenia Hurewicza,
danych normalnych i argumentu z twierdzenia~1.3
Walla. Nie wynika z samego ciągu dokładnego.

Ustawmy $T_x,T_y$ poprzecznie. Ich przecięcie
ma wymiar $(k+1)+(k+1)-(2k+1)=1$,
a więc składa się z okręgów i łuków.
Okręgi nie mają brzegu. Końce łuków
odpowiadają punktom przecięcia
$\partial T_x$ z $\partial T_y$ na $M$.
Po podniesieniu do nakrycia wyznaczonego
przez $\pi_1(Y)$ etykieta $g\in\pi_1(Y)$
jest stała wzdłuż jednego łuku, podczas
gdy orientacja jego dwóch końców daje
przeciwne znaki. Dla każdego $g$ suma
współczynników przy $g$ jest zatem zerem:
\[
 \lambda(\partial x,\partial y)=0.
\]
W szczególności $\widetilde\rho$ znika
na $\partial A$ i przechodzi do
\begin{equation}\label{eq:wall-57-dualnosc-22}
 \rho:C=P/\partial A\longrightarrow A^*,
 \qquad \rho([z])(x)=\lambda(\partial x,z).
\end{equation}
Jest to izomorfizm dualności względnej
z lematu~2.6 Walla, z dokładnością do
stałego znaku wynikającego z konwencji
orientacji. Zgodność dualności względnej,
operatora brzegu i formy przecięcia na $M$
utożsamia ocenę dualną z $\lambda(\partial x,z)$.
Ta zgodność jest osobnym wejściem
geometrycznym z dowodu lematu~5.7 Walla;
sam abstrakcyjny izomorfizm $C\cong A^*$
nie wyznacza wzoru na $\rho$.
Mnożenie przez $-1$ nie zmienia prostej
klasy tego izomorfizmu.

Tak samo ustawmy pojedynczą immersję
$T_x$ w położeniu ogólnym. Jej zbiór
podwójnych punktów ma wymiar jeden;
dla $k\geq2$ można uniknąć punktów
potrójnych przez pozycję ogólną.
Wzdłuż każdego łuku podwójnych punktów
ustalamy kolejność dwóch gałęzi.
Na jego końcach leżą samoprzecięcia
$\partial T_x$ o tej samej etykiecie
$g$ i przeciwnych znakach. Jeśli
zmienimy kolejność gałęzi, oba wkłady
przejdą na równoważne klasy w
$Q_{(-1)^k}(\Lambda)$.
Okręgi znów nie mają końców. Dlatego
\[
 \mu(\partial x)=0.
\]

\emph{Krok trzeci: lagranżjan i klasa $L$.}
Niech $L=\partial A\subset P$.
Z~\eqref{eq:wall-57-ciag-22} mamy $P/L\cong C$,
a~\eqref{eq:wall-57-dualnosc-22} mówi, że
\[
 P/L\xrightarrow{\;[z]\mapsto
           \lambda(-,z)|_L\;}L^*
\]
jest prostym izomorfizmem. Ponieważ
$\lambda|_{L\times L}=0$, jądro mapy
$P\to L^*$ równa się $L$; jest nim
z definicji $L^\perp$. Wraz z
$\mu|_L=0$ daje to lagranżjan kwadratowy.
Preferowana baza $A$ przechodzi na bazę
$L$; prostota dualności daje bazę $C$,
a bazowość~\eqref{eq:wall-57-ciag-22}
pozwala przedłużyć bazę $L$ do
preferowanej bazy $P$.
Lemat~\ref{lem:lagr-hip-Z-22} czyni
$(P,\lambda,\mu)$ formą hiperboliczną.

Teraz stosujemy ten rachunek do
$Y=X\times I$ i $M=\partial N=M_-\sqcup M_+$.
Mapy końców przygotowujemy do
$k$-spójności, mapę $N\to Y$ także,
a odejmowanie uchwytów usuwa
$K_k(N,\partial N)$. Po stabilizacji
preferowana stabilna baza $A$ ma
rzeczywistego reprezentanta.
W nieparzystym wymiarze $2k+1$
twierdzenie~1.4 Walla daje jednocześnie
$k$-spójność map $N\to Y$ i $\partial N\to X$
oraz $(k+1)$-spójność homologiczną mapy par.
W jego dowodzie ostatni warunek uzyskuje się
po zwykłych chirurgiach przez odejmowanie
uchwytów: generatory względnego jądra
przedstawia się obramowanymi dyskami
dochodzącymi do brzegu. Dla nowej mapy
$\Phi_0$ fragment ciągu dokładnego ma postać
\[
 \Lambda^r\longrightarrow H_{k+1}(\Phi)
 \longrightarrow H_{k+1}(\Phi_0)\longrightarrow0.
\]
Pierwsza strzałka jest surjekcją z wyboru
generatorów, więc ostatni moduł znika.
Dodatek do twierdzenia~1.4 pozwala wykonać
te zmiany po stronie $M_+$, zachowując $M_-$;
na brzegu odpowiadają one chirurgiom na
trywialnie obramowanych sferach $S^{k-1}$.
Lemat~5.5 Walla mówi, że każda taka
stabilizacja dodaje do formy $M_+$ standardową
płaszczyznę hiperboliczną. Ten krok zapewnia
też rzeczywiste reprezentanty potrzebnych
stabilnych baz; nie wynika z rachunku
lagranżjanu. Wówczas
\[
 P=K_k(\partial N)
 \cong K_k(M_+)\oplus K_k(M_-)^{-},
\]
gdzie minus oznacza przeciwną formę.
Hiperboliczność $P$ daje
$[K_k(M_+)]=[K_k(M_-)]$.
Tak otrzymujemy niezmienniczość
normalnej klasy kobordyzmu, z jawnym
użyciem twierdzeń~1.3 i~1.4 oraz
lematów~2.4, 2.6 i~5.5 Walla.
\end{proof}

\subsection*{Pełna wersja parzysta nad \texorpdfstring{$\Z[\pi]$}{Z[pi]}}

W tym ustępie $k\geq3$, $n=2k$, $\pi$ jest dowolną grupą,
$w:\pi\to\{\pm1\}$ charakterem orientacji celu, a
$\Lambda=\Z[\pi]$ ma inwolucję
$\overline g=w(g)g^{-1}$. Kładziemy
$\epsilon=(-1)^k$. Moduły są lewe; formę przecięcia
$\lambda:P\times P\to\Lambda$ zapisujemy liniowo w pierwszej
zmiennej, zatem
\[
 \lambda(ax,by)=a\lambda(x,y)\bar b,
 \qquad \lambda(y,x)=\epsilon\overline{\lambda(x,y)}.
\]
Współczynniki $\Lambda$ pochodzą zawsze od grupy podstawowej
\emph{celu}, także wtedy, gdy źródło jest niespójne.

Przez \emph{preferowaną bazę stabilną} modułu jądra rozumiemy
klasę baz uzyskaną z baz komórkowych kompleksu jądra przez
proste redukcje i stabilizacje. Dwie takie bazy mają tę samą
klasę, gdy odpowiednia zmiana bazy ma zerową klasę
w $\operatorname{Wh}(\pi)$. Zwykłe stwierdzenie, że moduł jest
wolny, nie zawiera tej informacji. Po dodaniu dostatecznie wielu
modułów wolnych preferowaną bazę stabilną można reprezentować
rzeczywistą bazą.

\begin{twierdzenie}[Twierdzenie 5.6 Walla w wymiarze parzystym $2k\geq6$]
\label{thm:Wall-parzyste-pelne-22}
Niech $(X,W)$ będzie spójną prostą parą Poincarégo o formalnym
wymiarze $2k\geq6$, $\nu\to X$ wiązką wektorową, a $M^{2k}$
zwartą gładką rozmaitością z brzegiem. Niech
$\phi:(M,\partial M)\to(X,W)$ będzie odwzorowaniem stopnia jeden,
którego ograniczenie $\partial M\to W$ jest prostą równoważnością
homotopijną. Wybierzmy stabilną trywializację
$F:\tau_M\oplus\phi^*\nu\cong M\times\R^N$ dla dostatecznie
dużego $N$.

Po chirurgiach normalnych relatywnie do $\partial M$ można
uczynić $\phi$ $k$-spójnym. Niech $\widetilde X$ będzie
uniwersalnym nakryciem $X$, a $\widetilde M=\phi^*\widetilde X$.
Środkowe jądro
\[
 K_k(M)=\ker\bigl(H_k(\widetilde M;\Z)
                \xrightarrow{\,\phi_*\,}H_k(\widetilde X;\Z)\bigr)
\]
z preferowaną bazą stabilną, formą przecięcia $\lambda$
i udoskonaleniem samoprzecięciowym $\mu$ wyznacza klasę
\[
 \theta(\phi,F)=[K_k(M),\lambda,\mu]
       \in L^s_{2k}(\Lambda,w).
\]
Nie zależy ona od wykonanych chirurgii ani od normalnego
kobordyzmu relatywnie do $\partial M$. Zachodzi
$\theta(\phi,F)=0$ wtedy i tylko wtedy, gdy tę mapę normalną
można przez taki kobordyzm zastąpić mapą, która jest prostą
równoważnością homotopijną. Przeszkoda jest naturalna dla
inkluzji w sensie opisanym na końcu dowodu.
\end{twierdzenie}

Wall definiuje tę klasę już dla $k\geq2$; równoważność
„$\theta=0$ wtedy i tylko wtedy, gdy można ukończyć chirurgię”
wymaga $k\geq3$. Symbol $L_{2k}(\pi)$ w jego §~5 oznacza tutaj
grupę z dekoracją $s$: istotne są proste bazy, prosta dualność
i prosta równoważność końcowa.

\noindent\textbf{Wejścia zewnętrzne.}
Poniższy dowód przyjmuje następujące dokładne wyniki
\href{https://www.maths.ed.ac.uk/~v1ranick/books/scm.pdf}
{Walla, \emph{Surgery on Compact Manifolds}, §§~1--2 i~5}:
twierdzenie~1.2 o chirurgiach poniżej środka relatywnie do
brzegu; twierdzenie~1.4 o przygotowaniu kobordyzmu przez
odejmowanie uchwytów; lematy~2.3, 2.4 i~2.6 o preferowanych
bazach, bazowanych ciągach dokładnych i prostej dualności;
lemat~5.1 o stabilnej bazie i prostej formie środkowego jądra;
twierdzenie~5.2 o przecięciach oraz o zanurzeniu obramowanej
sfery reprezentującej klasę z $\mu=0$ przy $k\geq3$;
twierdzenie~1.1 o klasach regularnej homotopii
obramowanych immersji;
lemat~5.7 o prostym lagranżjanie na brzegu przygotowanego
kobordyzmu. Twierdzenie~1.1 pozwala wykonać chirurgię
na otrzymanej zanurzonej sferze. Są to twierdzenia
geometryczne, a nie wnioski z poniższych dwóch ciągów
dokładnych. W szczególności twierdzenie~5.2 używa
obramowanego triku Whitneya w środku wymiaru.

\begin{lemat}[Sprowadzenie chirurgii na klasie zerowej do dysku]
\label{lem:wall-zero-do-dysku-22}
W hipotezach Twierdzenia~\ref{thm:Wall-parzyste-pelne-22}
niech $k\geq3$, a
$e:S^{k-1}\times D^{k+1}\hookrightarrow\operatorname{int}M$
będzie obramowanym zanurzeniem dopuszczalnym dla
danych normalnych $(\phi,F)$. Jeżeli klasa
wyznaczona przez rdzeń $e$ wraz z wybraną
homotopią zerującą jego obraz w $X$ jest zerem
w $\pi_k(\phi)$, to $e$ można przeprowadzić
izotopią zachowującą obramowanie na model
standardowy w pewnym dysku
$D^{2k}\subset\operatorname{int}M$.
\end{lemat}
\begin{proof}
W dysku wybierzmy standardową sferę
$S^{k-1}=\partial D^k$ z produktowym otoczeniem.
Kontrakcja dysku daje homotopię zerującą
obraz tej sfery w $X$; trywializacja $F$
na dysku wyznacza zgodne obramowanie.
Jest to model klasy zerowej.
Twierdzenie~1.1
\href{https://www.maths.ed.ac.uk/~v1ranick/books/scm.pdf}
{Walla, §1} przypisuje klasie względnej
\emph{jedną klasę regularnej homotopii
obramowanych immersji}. Dlatego $e$ i model
w dysku łączy taka homotopia. Sama zerowość
obrazu sfery w $X$ nie wystarczyłaby do
tego wniosku: potrzebne są wybrana homotopia
zerująca i zgodność ram z danymi normalnymi.

Zastosujmy parametryczne położenie ogólne
do śladu homotopii rdzeni, relatywnie do
jej końców. W rodzinie jednoparametrowej
oczekiwany wymiar zbioru punktów podwójnych
wynosi
\[
 2(k-1)+1-2k=-1.
\]
Możemy więc wybrać ślad bez takich punktów;
każda sfera pośrednia jest zanurzona, a rodzina
rdzeni jest izotopią we wnętrzu $M$.
Rzeczywiste obramowania również dają się
przenieść wzdłuż tej izotopii: ich ranga to
$k+1$, a stabilizacja
$\pi_{k-1}(O(k+1))\to\pi_{k-1}(O)$
jest izomorfizmem, gdyż kolejne ilorazy
$O(q+1)/O(q)$ są sferami $S^q$.
Zgodność stabilnych ram z twierdzenia~1.1
oznacza zatem zgodność ram rzeczywistych.
Zwężamy włókna obu otoczeń rurkowych i
przedłużamy obramowaną izotopię rdzeni
na otoczenia. Jej podparcie omija $\partial M$.
\end{proof}

\begin{lemat}[Co dokładnie dodaje lokalna stabilizacja]
\label{lem:wall-lokalna-stabilizacja-22}
W hipotezach Twierdzenia~\ref{thm:Wall-parzyste-pelne-22}
niech $\phi$ będzie już $k$-spójne. Chirurgia na
dowolnej dopuszczalnej obramowanej $(k-1)$-sferze
reprezentującej zero w $\pi_k(\phi)$ zmienia środkowe
jądro i jego formę według wzoru
\[
 (K_k(M'),\lambda',\mu')
 \cong (K_k(M),\lambda,\mu)
       \perp H_{(-1)^k}(\Lambda),
\]
gdzie nowa płaszczyzna ma preferowaną bazę
$e,f$ z $\lambda(e,f)=1$ i
$\mu(e)=\mu(f)=0$.
\end{lemat}
\begin{proof}
Lemat~\ref{lem:wall-zero-do-dysku-22} sprowadza
tę chirurgię przez obramowaną izotopię do
standardowej chirurgii w dysku.
Lokalny rachunek wykonujemy najpierw w
$S^{2k}=\partial(D^k\times D^{k+1})$.
Jego rozkład brzegowy zawiera
$S^{k-1}\times D^{k+1}$ oraz
$D^k\times S^k$. Zastąpienie pierwszej
części przez drugą w chirurgii daje
dwie kopie $D^k\times S^k$ sklejone po
$S^{k-1}\times S^k$, czyli
$S^k\times S^k$. Dlatego chirurgia w
dysku rozmaitości $M$ daje
$M'\cong M\#(S^k\times S^k)$.
Mapę do $X$ na nowym składniku można
wybrać stałą. W nakryciu indukowanym
z $\widetilde X$ każda kopia tego
składnika dostarcza dwie klasy
$S^k\times\{*\}$ i $\{*\}\times S^k$.
Ich przesunięcia przez elementy $\pi$
tworzą wolny składnik
$\Lambda e\oplus\Lambda f$ w jądrze;
ciąg Mayera--Vietorisa dla sumy
spójnej pokazuje, że nie powstają
relacje ze starym $K_k(M)$.

Obie sfery czynnikowe przecinają się
w jednym punkcie. Ich obramowania pochodzą
z produktowych uchwytów i są zgodne
z danymi normalnymi śladu. Tak obramowane
sfery są zanurzone i nie mają punktów podwójnych.
Twierdzenie~5.2 Walla
o obramowanych immersjach daje zatem
$\lambda(e,f)=1$ oraz $\mu(e)=\mu(f)=0$.
Leżą w dysku
rozłącznym od reprezentantów starego
jądra, więc suma jest ortogonalna.
Nowe generatory pochodzą od
standardowych uchwytów lokalnej
chirurgii; ich macierze przy
kontrakcjach mają jednostki $1$.
Zatem baza $(e,f)$ należy do
preferowanej prostej klasy baz.
\end{proof}

Oba powyższe lematy razem wyprowadzają
lemat~5.5 Walla w używanym tu zakresie,
przy jawnym zastosowaniu jego twierdzenia~1.1.

\begin{proof}[Dowód z wyszczególnieniem wejść Walla]
\textbf{1. Redukcja i właściwy obiekt algebraiczny.}
Twierdzenie~1.2 Walla pozwala przez skończenie wiele
chirurgii relatywnie do brzegu doprowadzić $\phi$ do
$k$-spójności. W tej sytuacji jądra nad $\Lambda$ znikają
poza stopniem $k$. Lemat~5.1 identyfikuje
$K_k(M)$ z odpowiednią względną grupą homotopii i nadaje
mu preferowaną bazę stabilną. Prosta dualność pary
Poincarégo daje prosty izomorfizm $K_k(M)\to K_k(M)^*$.
Nie wynika to z samej $k$-spójności: w zastosowaniu
lematu~2.3(c) Walla potrzebne są również skończoność
kompleksów i zanik pozostałych homologii kompleksu
jądra, a na odpowiednim krańcu zanik z dowolnymi
współczynnikami $\Lambda$-modułowymi. Zapewniają je
prosta dualność Poincarégo, warunek stopnia jeden
i prosta równoważność na brzegu, jak w dowodzie
lematu~5.1 Walla.
Twierdzenie~5.2 konstruuje na tym module formę
$\epsilon$-kwadratową $(\lambda,\mu)$.
To wyjaśnia, dlaczego jej klasę liczymy w $L^s_{2k}$,
a nie wyłącznie w grupie zwykłych form na abstrakcyjnych
modułach.

\textbf{2. Stabilizacja i baza hiperboliczna.}
Jeśli $\theta(\phi,F)=0$, definicja zerowej klasy
w parzystej grupie chirurgicznej mówi, że po dodaniu
pewnej liczby standardowych płaszczyzn hiperbolicznych
otrzymamy prostą formę hiperboliczną. Algebraiczne
przejście od prostego lagranżjanu do preferowanej bazy
hiperbolicznej jest rozpisane w
Lemacie~\ref{lem:lagr-hip-Z-22}; jego hipoteza o prostocie
$P/L\to L^*$ jest potrzebna właśnie dla torsji.

Każdą algebraicznie dodaną płaszczyznę realizuje
chirurgia z
Lematu~\ref{lem:wall-lokalna-stabilizacja-22}.
Możemy więc, po
normalnym kobordyzmie relatywnie do brzegu, założyć
\[
 K_k(M)=\bigoplus_{i=1}^{r}(\Lambda e_i\oplus\Lambda f_i),
 \quad \lambda(e_i,f_j)=\delta_{ij},\quad
 \lambda(e_i,e_j)=\lambda(f_i,f_j)=0,
 \quad \mu(e_i)=\mu(f_i)=0.
\tag{H}
\]
Baza $(e_1,f_1,\ldots,e_r,f_r)$ jest preferowana,
nie tylko dowolna.

\textbf{3. Jedna chirurgia środkowa i jej dwa ciągi.}
Lemat~5.1 i twierdzenie~5.2 przedstawiają $e_r$
przez obramowaną immersję $S^k\looparrowright M$.
Równość $\mu(e_r)=0$ pozwala, gdy $k\geq3$,
zastąpić ją obramowanym zanurzeniem
$S^k\hookrightarrow M$ reprezentującym tę samą klasę
jądra. Chirurgia nad $X$ według twierdzenia~1.1
ma ślad $T$ i nowy brzeg $M'$. Kompleks pary
$(T,M)$ ma jeden generator w stopniu $k+1$,
a kompleks pary $(T,M')$ jeden generator w stopniu $k$;
w obu przypadkach współczynnikiem jest $\Lambda$.

Pierwszy ciąg pary, z uwzględnieniem jądra mapy
do celu, zawiera strzałkę $\Lambda\to K_k(M)$,
która wysyła $1$ na $e_r$. Jest injektywna,
ponieważ $e_r$ należy do bazy (H). Otrzymujemy
\begin{equation}\label{eq:Wall-parz-pierwszy-22}
 0\longrightarrow\Lambda e_r\longrightarrow K_k(M)
  \longrightarrow K_k(T)\longrightarrow0.
\end{equation}
Od drugiego końca uchwytu dostajemy
\begin{equation}\label{eq:Wall-parz-drugi-22}
 0\longrightarrow K_k(M')\longrightarrow K_k(T)
  \xrightarrow{\,\ell_{e_r}\,}\Lambda
  \longrightarrow K_{k-1}(M')\longrightarrow0.
\end{equation}
Strzałka $\ell_{e_r}$ mierzy przecięcie z osią
dołączonego uchwytu. Przy naszej konwencji lewych
modułów jest indukowana przez
\[
 \ell_{e_r}(z)=\epsilon\lambda(z,e_r)
 \qquad(z\in K_k(M)).
\]
Jest liniowa w $z$, znika na $e_r$ i
$\ell_{e_r}(f_r)=
\epsilon\lambda(f_r,e_r)=\epsilon^2=1$.
Ta ostatnia równość daje rozszczepienie strzałki
w~\eqref{eq:Wall-parz-drugi-22}, a więc
$K_{k-1}(M')=0$. Jej jądro w~$K_k(M)$ to
$e_r^\perp=\{z:\lambda(z,e_r)=0\}$, skąd
\begin{equation}\label{eq:Wall-parz-iloraz-22}
 K_k(M')\cong e_r^\perp/\Lambda e_r
   \cong\bigoplus_{i<r}(\Lambda e_i\oplus\Lambda f_i).
\end{equation}
Ponieważ $\mu(e_r)=0$ i forma znika między
$e_r$ a $e_r^\perp$, udoskonalenie $\mu$
przechodzi na ten iloraz. Jest to dokładnie forma
hiperboliczna na pozostałych parach, z ich
preferowaną bazą.

Ostatnie zdanie wymaga kontroli \emph{prostej}
klasy bazy. Na płaszczyźnie $e_r,f_r$ obie
strzałki tworzą bazowany kompleks
\[
 0\longrightarrow\Lambda e_r
 \xrightarrow{\,e_r\mapsto e_r\,}
 \Lambda e_r\oplus\Lambda f_r
 \xrightarrow{\,e_r\mapsto0,\;f_r\mapsto1\,}
 \Lambda\longrightarrow0.
\]
Jego macierze zawierają jednostkę $1$, więc
elementarna kontrakcja ma zerową torsję
Whiteheada. Geometryczne bazy komórek uchwytowych
i bazowany wniosek z lematu~2.4 Walla pozwalają
przenieść ten rachunek na preferowane stabilne
bazy kompleksów jądra.
Dlatego izomorfizm~\eqref{eq:Wall-parz-iloraz-22}
zachowuje preferowaną klasę bazy; zwykły
izomorfizm modułów nie wystarczałby do tego wniosku.

Powtarzając krok dla $e_{r-1},\ldots,e_1$,
usuwamy całe $K_k$ i nie tworzymy niższego jądra.
Końcowa mapa jest $k$-spójna i ma zerowe homologie
kompleksu jądra nad $\Lambda$, zatem jest
równoważnością homotopijną przez twierdzenie
Whiteheada z lokalnymi współczynnikami.
Preferowane bazy i powyższe elementarne kontrakcje
dają zerową torsję jej stożka, czyli równoważność
jest \emph{prosta}. To kończy dostateczność.

\textbf{4. Konieczność i niezależność od wyborów.}
Jeśli końcowa mapa jest prostą równoważnością
homotopijną, jej bazowany kompleks jądra jest
prosto ściągalny. Jego klasa w $L^s_{2k}$ wynosi zero.

Dla niezmienniczości porównajmy dwie przygotowane
$k$-spójne mapy normalne $M_-\to X$ i $M_+\to X$
połączone normalnym kobordyzmem relatywnie do brzegu.
Jego celem jest para $(Y,X_0)$, gdzie
$Y=X\times I$ oraz
$X_0=X\times\{0,1\}\cup W\times I$.
Walla twierdzenia~1.2 i~1.4 pozwalają tak przygotować
mapę śladu $N^{2k+1}\to Y$, aby zarówno ona,
jak i mapa całego brzegu $\partial N\to X_0$
były $k$-spójne oraz aby $K_k(N,\partial N)=0$;
konieczne
zmiany na końcach dodają tylko płaszczyzny
hiperboliczne z lematu~5.5. Można następnie
stabilizować tak, by preferowana stabilna baza
$A=K_{k+1}(N,\partial N)$ była reprezentowana
przez rzeczywistą bazę. Wszystkie te operacje
są geometryczną treścią wymienionych twierdzeń,
nie konsekwencją samej algebraicznej klasy $\theta$.

Zapiszmy $P=K_k(\partial N)$ i $C=K_k(N)$;
wszystkie grupy liczymy nad
$\Z[\pi_1(X\times I)]=\Lambda$. Ponieważ
mapa na bocznym brzegu jest prostą równoważnością,
\[
 (P,\lambda,\mu)
  \cong (K_k(M_+),\lambda_+,\mu_+)
       \oplus(K_k(M_-),-\lambda_-,-\mu_-).
\]
W tej sytuacji bazowany ciąg dokładny ma postać
\begin{equation}\label{eq:Wall-parz-lagr-ciag-22}
 0\longrightarrow A\xrightarrow{\,\partial\,}P
   \longrightarrow C\longrightarrow0.
\end{equation}
\emph{Bazowanie} tego ciągu jest wnioskiem z~2.4
Walla, a prosta dualność $C\cong A^*$ z~2.6.
Diagram dualności Poincarégo utożsamia powstałe
$P/\partial A\cong C\cong A^*$ z odwzorowaniem
$[p]\mapsto(a\mapsto\lambda(\partial a,p))$;
stąd jego prostota jest właśnie prostotą
dualności formy na ilorazie.
Lemat~5.7 Walla mówi więcej: $\partial A$ jest
\emph{prostym lagranżjanem} formy na~$P$,
czyli $\lambda|_{\partial A}=0$,
$\mu|_{\partial A}=0$, jego preferowana baza
przedłuża się do preferowanej bazy $P$,
a indukowana dualność $P/\partial A\to A^*$
jest prosta. Jego dokładne hipotezy to: mapa normalna
stopnia jeden $(N,\partial N)\to(Y,X_0)$ ze zwartej
rozmaitości $N$ do prostej pary Poincarégo ze
spójnym $Y$, z danymi $\nu,F$ jak wyżej,
$\dim N=2k+1\geq5$, $k$-spójność map
$N\to Y$ i $\partial N\to X_0$,
$K_k(N,\partial N)=0$ oraz rzeczywista baza
reprezentująca preferowaną bazę stabilną
$K_{k+1}(N,\partial N)$. W zastosowaniu
$Y=X\times I$; przy ogólnej parze $(Y,X_0)$
trzeba używać $\Z[\pi_1(Y)]$, nie
$\Z[\pi_1(X_0)]$.

Lemat~5.3 Walla, czyli algebraiczny
Lemat~\ref{lem:lagr-hip-Z-22}, czyni formę
na~$P$ hiperboliczną. Z~\eqref{eq:Wall-parz-lagr-ciag-22}
i znaku orientacji obu końców wynika zatem
$[K_k(M_+)]=[K_k(M_-)]$ w $L^s_{2k}(\Lambda,w)$.
To dowodzi niezależności od kobordyzmu oraz,
biorąc kobordyzmy realizujące różne przygotowania,
od wyboru chirurgii poniżej środka.

\textbf{5. Naturalność dla inkluzji.}
Niech przy $W\subset X$ do pary $(X,W)$ będzie
doklejona triada przez prostą równoważność
homotopijną triad, jak w ostatnim akapicie dowodu
twierdzenia~5.6 Walla. Dla $X\hookrightarrow X'$
otrzymujemy homomorfizm grup i pierścieni
$\rho:\Z[\pi_1X]\to\Z[\pi_1X']$ zgodny
z charakterami orientacji. Nowa część triady
ma prosto ściągalny kompleks jądra, więc po
zmianie współczynników jądro starej części,
jego preferowana baza oraz liczby przecięć
i samoprzecięć przechodzą odpowiednio na
$\Z[\pi_1X']\otimes_{\Lambda}K_k(M)$,
$\rho(\lambda)$ i $\rho(\mu)$. Stąd
\[
 \theta(\phi\cup\phi',F\cup F')
       =\rho_*\theta(\phi,F)
       \quad\text{w }L^s_{2k}(\Z[\pi_1X'],w').
\]
Nie jest to stwierdzenie o dowolnym odwzorowaniu
celów: chodzi o tę konkretną inkluzję z dołączoną
prostą równoważnością triad.
\end{proof}

\noindent\textbf{Zakres dowodu.}
Rachunek~\eqref{eq:Wall-parz-pierwszy-22}--
\eqref{eq:Wall-parz-iloraz-22}, usunięcie pary
hiperbolicznej oraz jej elementarna torsja są
wyprowadzone powyżej. Istnienie wymaganych
obramowanych zanurzeń, przygotowanie kobordyzmu,
preferowane bazy i prosta dualność są użytymi
z dokładnymi hipotezami wynikami Walla.
Fragment zamyka \emph{parzysty} przypadek $n\geq6$;
przypadek nieparzysty używa formacji i odrębnych
twierdzeń jego §~6.

\begin{figure}[htbp]
\centering
\begin{tikzpicture}[>=Latex,font=\small,
 box/.style={draw=manifoldblue!75!black,rounded corners=3pt,
 fill=manifoldblue!9,align=center,minimum height=1.05cm,
 text width=2.4cm}]
 \node[box] (a) at (0,0) {normalna mapa\\$M\to X$};
 \node[box] (b) at (3.25,0)
   {chirurgie poniżej\\środka};
 \node[box] (c) at (6.50,0)
   {forma jądra\\$(K,\lambda,\mu)$};
 \node[box] (d) at (9.75,0)
   {lagranżjan i\\chirurgie środkowe};
 \draw[->,thick,manifoldblue!75!black]
   (a)--(b);
 \draw[->,thick,manifoldblue!75!black]
   (b)--(c);
 \draw[->,thick,manifoldblue!75!black]
   (c)--(d);
 \draw[->,thick,manifoldgold!85!black]
   (d.south) -- ++(0,-1.30) -| (a.south);
 \node[fill=white,font=\scriptsize,align=center]
   at (4.85,-1.30) {zanik jądra daje równoważność};
\end{tikzpicture}
\caption{Kolejne etapy twierdzenia Walla w parzystym wymiarze.
Przejście od algebraicznego lagranżjanu do zanurzonych
sfer wymaga lematów geometrycznych Walla.}
\label{fig:mechanizm-Wall-22}
\end{figure}

\subsection*{Dlaczego wymiar nieparzysty wymaga formacji}
Niech $n=2k+1\geq5$ i niech mapa normalna
została przygotowana chirurgiami poniżej
środka. Po stabilizacji wybieramy rozłączne
obramowane reprezentanty generatorów
$K_k(M)$, z otoczeniem
$U=\coprod(S^k\times D^{k+1})$.
To wybranie, wraz z kontrolą prostych baz,
jest geometrycznym krokiem z §6 Walla.
Połóżmy $C=\overline{M\setminus\operatorname{int}U}$.
Moduł $H=K_k(\partial U)$ ma formę
hiperboliczną: na każdej składowej brzegu
$S^k\times S^k$ dwie sfery czynnikowe
przecinają się raz. Po wyborze dróg bazowych
dają bazę $e_1,f_1,\ldots,e_r,f_r$ z
$\lambda(e_i,f_j)=\delta_{ij}$ oraz
$\mu(e_i)=\mu(f_i)=0$.
W obrazie odwzorowań brzegu dostajemy
dwa lagranżjany
\[
 L_0=\operatorname{im}\bigl(K_{k+1}(U,\partial U)
                  \to H\bigr),\qquad
 L_1=\operatorname{im}\bigl(K_{k+1}(C,\partial U)
                  \to H\bigr).
\]
Formacja jest trójką $(H;L_0,L_1)$.
Pierwszy lagranżjan jest jawnie rozpięty przez
jedną rodzinę sfer czynnikowych. Dla drugiego
samo zapisanie go jako obrazu strzałki
\emph{nie} dowodzi warunku $L_1=L_1^\perp$.
Wall najpierw przesuwa obraz $U$ do dysku w~$X$,
przygotowuje parę $(C,\partial U)$ oraz jej
bazy, a następnie stosuje lemat~5.7 do tej pary.
Prostota bazy $L_1$ korzysta tu z lematów~2.4
i~2.6 oraz ze stabilizacji przez dołożenie
reprezentantów generatorów. Ta część konstrukcji
pozostaje wejściem geometrycznym z~§6 Walla.
Jej użycie nie oznacza, że dwa sąsiednie
jądra \emph{są} lagranżjanami; powstają
z ich wzajemnego położenia.

Algebraiczny ciąg
\begin{equation}\label{eq:formacja-ciag-Wall-22}
 0\longrightarrow L_0\cap L_1
 \longrightarrow L_1
 \longrightarrow H/L_0
 \longrightarrow H/(L_0+L_1)
 \longrightarrow0
\end{equation}
jest dokładny: środkowa strzałka wysyła
$z$ na $[z]$, więc jej jądrem jest
$L_0\cap L_1$, a jej obrazem
$(L_0+L_1)/L_0$.
Diagram Mayera--Vietorisa dla $M=U\cup C$
identyfikuje skrajne człony z
$K_{k+1}(M)$ i $K_k(M)$, odpowiednio.
Jeśli $L_0\oplus L_1=H$, oba te jądra
znikają; po zaniku niższych jąder i
przy $\pi_1=1$ twierdzenie Whiteheada
daje równoważność prostą.

Można obliczać oba skrajne człony jedną
macierzą. Po wyborze baz $H=E\oplus E^*$,
$L_0=E$, a prosta izometria $\alpha:H\to H$
przenosi $E$ na $L_1$; jej istnienie dla
prostych lagranżjanów to wniosek~5.3.1 Walla.
Zdefiniujmy
\[
 T=\operatorname{pr}_{E^*}\circ\alpha|_E:
 E\longrightarrow E^*.
\]
Ponieważ iloraz $H/L_0$ jest kanonicznie
$E^*$, środkowa strzałka
\eqref{eq:formacja-ciag-Wall-22}, po
utożsamieniu $L_1$ z $E$ przez $\alpha$,
jest dokładnie $T$. Zatem
\[
 K_{k+1}(M)\cong\ker T,\qquad
 K_k(M)\cong\operatorname{coker}T.
\]
To jest rachunek jąder dla \emph{wybranej}
formacji. Niezależność jej stabilnej klasy
od $U$ oraz wpływ chirurgii na $\alpha$
wymagają osobnych ruchów z~§6 Walla.

Na przykład dla $\epsilon=-1$ weźmy
$H=\Z e\oplus\Z f$ z
$\lambda(e,f)=1$ oraz
$\mu(ae+bf)=ab\pmod2$.
Podmoduły $L_0=\Z e$ i
$L_1=\Z(e+2f)$ są lagranżjanami:
$\mu(e+2f)=0$, a wektor $e+2f$
jest prymitywny. Tutaj $T$ jest
mnożeniem przez $2$: strzałka
$L_1\to H/L_0\cong\Z f$ z
\eqref{eq:formacja-ciag-Wall-22} jest
mnożeniem przez $2$. Jej cokernel to
$\Z/2$, zatem sam fakt istnienia dwóch
lagranżjanów nie oznacza, że obecne
jądro już znikło. Mimo to \emph{klasa tej
formacji jest zerowa}: weźmy jednoelementowy
moduł $K=\Z e$ z parzystą formą
$\gamma(e,e)=2$ i jej udoskonaleniem
kwadratowym $\mu_\gamma(e)=1$.
W hiperbolicznym $K\oplus K^*$ wykres
odwzorowania $\gamma:K\to K^*$ jest
$\Z(e+2f)$. Nasza para $(L_0,L_1)$ jest więc
\emph{formacją brzegową} $\partial\gamma$,
która z definicji reprezentuje zero w grupie
$L^s_{2k+1}(\Z)$. Ten rachunek pokazuje
konkretnie różnicę między zerem przeszkody
a zanikiem jąder przed wykonaniem dalszych
chirurgii.

Przesunięcie uchwytu ma czytelny rachunek
w $H_\epsilon\oplus H_\epsilon$:
$e_1\mapsto e_1+e_2$ oraz
$f_2\mapsto f_2-f_1$, przy pozostałych
wektorach stałych, zachowuje wszystkie
parowania i $\mu$. Pełna relacja formacji
obejmuje ponadto stabilizację i formacje
brzegowe. Dlatego zerowa \emph{klasa} w
$L^s_{2k+1}(\Lambda,w)$ nie wymaga, aby
wybrane już teraz $L_0$ i $L_1$ były
dopełniające. Realizacja dozwolonych
ruchów przez chirurgie oraz niezależność
wyboru $U$ to treść
\href{https://webhomes.maths.ed.ac.uk/~v1ranick/books/scm.pdf}
{Walla, §6, twierdzenia 6.4 i 6.6};
algebraiczną relację formacji wyjaśnia
\href{https://www.maths.ed.ac.uk/~v1ranick/papers/ranicki.pdf}
{A.\ Ranicki, \emph{An Introduction to Algebraic Surgery}, §9}.

\subsection*{Pełne kryterium chirurgiczne w wymiarze nieparzystym}

Niech $m=2k+1\geq5$, niech
$\pi=\pi_1(X)$ i niech $w:\pi\to\{\pm1\}$ będzie
homomorfizmem orientacji. Na pierścieniu
$\Lambda=\Z[\pi]$ stosujemy antyinwolucję
\[
 \overline{\sum_g n_g g}=\sum_g n_g w(g)g^{-1},
 \qquad \epsilon=(-1)^k.
\]
W tym podrozdziale wszystkie moduły są lewymi
$\Lambda$-modułami. Oznaczenia formy
$\epsilon$-kwadratowej i lagranżjanu są takie jak
w~Definicji~\ref{def:forma-kwadratowa-22} i
Definicji~\ref{def:hiperboliczna-22}. Dla czytelnika
kluczowa różnica względem wymiaru parzystego jest
następująca: dwa środkowe jądra jednej mapy
powstają z \emph{wzajemnego położenia dwóch
lagranżjanów} w jednym module hiperbolicznym.

\paragraph{Grupa algebraiczna.}
Niech $E=\Lambda^r$ z ustaloną bazą i niech
$E^*=\operatorname{Hom}_\Lambda(E,\Lambda)$, z
lewym działaniem $(a\varphi)(x)=\varphi(x)\bar a$.
Na $H_\epsilon(E)=E\oplus E^*$ połóżmy
\[
 \lambda\bigl((x,\varphi),(y,\psi)\bigr)
   =\psi(x)+\epsilon\overline{\varphi(y)},
 \qquad
 \mu(x,\varphi)=[\varphi(x)]\in
 Q_\epsilon(\Lambda)
 =\Lambda/\{a-\epsilon\bar a:a\in\Lambda\}.
\]
Iloraz w ostatnim wzorze jest ilorazem grup
addytywnych. Dla bazy $e_i$ modułu $E$ i jej
bazy dualnej $f_i$ mamy
$\lambda(e_i,f_j)=\delta_{ij}$ oraz
$\mu(e_i)=\mu(f_i)=0$.

Piszemy $SU_r(\Lambda)$ na grupę automorfizmów
$(H_\epsilon(E),\lambda,\mu)$ zachowujących
wyróżnioną prostą klasę baz. Słowo „prosty”
oznacza tutaj zerową torsję Whiteheada
automorfizmu modułu; skrót $SU$ nie określa
wyznacznika. Wyróżniamy podgrupy
\[
 TU_r(\Lambda)=\{A\in SU_r(\Lambda):
             A(E)=E,\ A|_E\text{ jest prosty}\},
 \qquad
 UU_r(\Lambda)=\{A\in TU_r(\Lambda):A|_E=1\}.
\]
Ograniczenie do $E$ daje rozszczepiony ciąg
\[
 1\longrightarrow UU_r(\Lambda)
 \longrightarrow TU_r(\Lambda)
 \longrightarrow\operatorname{Aut}^{s}_\Lambda(E)
 \longrightarrow1.
\]
Rozszczepienie wysyła $B$ na automorfizm
hiperboliczny
$H(B)(x,\varphi)=(Bx,\varphi\circ B^{-1})$.
W bazie $e_i,f_i$ element $UU_r$ ma postać
\[
 u_C(x,\varphi)=(x+C\varphi,\varphi),
 \qquad C=D-\epsilon D^*,
\]
gdzie $D^*$ jest macierzą sprzężoną i
transponowaną. Sam warunek zachowania
$\lambda$ daje tylko $C+\epsilon C^*=0$;
postać $D-\epsilon D^*$ uwzględnia dodatkowo
zachowanie $\mu$ na przekątnej.

Niech $\sigma\in SU_1(\Lambda)$ zamienia dwa
standardowe lagranżjany według wzorów
$\sigma(e)=f$ i $\sigma(f)=\epsilon e$; zatem
$\sigma^2=\epsilon\,1$. Grupy $SU(\Lambda)$,
$TU(\Lambda)$ otrzymujemy przez stabilizację
sumą ortogonalną z identycznością na jednej
płaszczyźnie hiperbolicznej. Niech
\[
 RU(\Lambda)=\langle TU(\Lambda),\sigma\rangle
       \subset SU(\Lambda).
\]
Nie zakładamy tu z góry normalności $RU$.
\href{https://www.maths.ed.ac.uk/~v1ranick/books/scm.pdf}
{Lemat 6.2 i twierdzenie 6.3 Walla, §6}
dają kolejno, po stabilizacji,
$A\oplus A^{-1}\in RU(\Lambda)$ dla każdego
$A\in SU(\Lambda)$ oraz
$[SU(\Lambda),SU(\Lambda)]\subset RU(\Lambda)$.
Stąd $RU$ jest normalna, a
\[
 L^s_{2k+1}(\Lambda,w)
       :=SU(\Lambda)/RU(\Lambda)
\]
jest grupą abelową. W dowodzie Walla lemat~6.2
jest jawną tożsamością blokowych macierzy
hiperbolicznych; twierdzenie~6.3 używa jej do
wykazania normalności. Te dwa wyniki algebraiczne
są tutaj cytowanymi lematami zewnętrznymi.
W poniższym zapisie automorfizmy składamy jak
odwzorowania: $(AB)(x)=A(Bx)$. Wall zapisuje
działanie macierzy na bazach z prawej, dlatego
kolejność czynników w jego wzorach dla kosetów
jest odwrotna.

\begin{twierdzenie}[Wall, twierdzenie 6.4:
                    przypadek nieparzysty]
\label{thm:wall-odd-22}
Niech $(X,W)$ będzie spójną prostą parą Poincarégo
w formalnym wymiarze $m=2k+1\geq5$, a $\nu$ wiązką
stabilną nad $X$. Niech $M^m$ będzie zwartą gładką
rozmaitością z brzegiem, zaś
$f:(M,\partial M)\to(X,W)$ odwzorowaniem stopnia
jeden, którego ograniczenie
$\partial M\to W$ jest prostą równoważnością
homotopijną. Niech $F$ będzie stabilną
trywializacją $TM\oplus f^*\nu$.
Istnieje klasa
\[
 \theta(f,F)\in L^s_m(\Z[\pi_1(X)],w),
\]
która zależy tylko od klasy kobordyzmu normalnego
$(f,F)$ relatywnie do brzegu. Klasa ta jest zerowa
wtedy i tylko wtedy, gdy $(f,F)$ jest normalnie
kobordantne relatywnie do brzegu z prostą
równoważnością homotopijną. Przeszkoda jest
naturalna przy inkluzji par Poincarégo, dla której
dane brzegowe przedłużają się przez prostą
równoważność homotopijną.
\end{twierdzenie}

\begin{lemat}[Rozłączne sfery w nieparzystym wymiarze]
\label{lem:wall-odd-rozlaczne-sfery-22}
W hipotezach Twierdzenia~\ref{thm:wall-odd-22}
załóżmy, że $f$ jest już $k$-spójne.
Każda skończona rodzina klas
$x_1,\ldots,x_r\in\pi_{k+1}(f)$, wraz
z drogami bazowymi i homotopiami zerującymi
ich obrazy w $X$, ma reprezentantów w postaci
parami rozłącznych obramowanych zanurzeń
\[
 e_i:S^k\times D^{k+1}\hookrightarrow
       \operatorname{int}M.
\]
Dla $U=\bigcup_i e_i(S^k\times D^{k+1})$
można po homotopii $f$ relatywnie do $\partial M$
przyjąć $f(U)=*$, zachowując dane normalne.
\end{lemat}
\begin{proof}
Klasę $x_i$ przedstawiamy parą
$u_i:S^k\to M$ i $H_i:D^{k+1}\to X$,
gdzie $H_i|_{S^k}=f\circ u_i$.
Przez kołnierz brzegu przesuwamy obrazy
$u_i$ do wnętrza $M$.
\href{https://www.maths.ed.ac.uk/~v1ranick/books/scm.pdf}
{Twierdzenie 1.1 Walla} wyznacza z $H_i$
i stabilnych danych $F$ klasę regularnej
homotopii obramowanych immersji pogrubień
$S^k\times D^{k+1}\looparrowright M$.
Wybieramy takie immersje i nadal oznaczamy
ich rdzenie przez $u_i$.
Homotopie $H_i$ przenosimy wzdłuż homotopii
rdzeni, tak by nadal zachodziło
$H_i|_{S^k}=f\circ u_i$ i klasy $x_i$
pozostały te same.
W szczególności obramowanie ma być zgodne
z \emph{wybraną} homotopią zerującą, a nie
tylko z klasą samego odwzorowania $u_i$.

Możliwość wyboru rzeczywistych ram widać
również bezpośrednio. Wiązka $H_i^*\nu$
nad dyskiem jest trywialna, zatem $F$
i $TS^k\oplus\R\cong\R^{k+1}$ dają
stabilne obramowanie wiązki normalnej
rdzenia $u_i$, rangi $k+1$.
Inkluzja $O(k+1)\to O$ indukuje izomorfizm
na $\pi_{k-1}$, więc stabilnie trywialna
wiązka nad $S^k$ jest rzeczywiście trywialna.
Na $\pi_k$ ta inkluzja jest surjekcją;
skręcając rzeczywiste obramowanie, można
więc dopasować jego klasę stabilną do $F$
i $H_i$. Oba stwierdzenia wynikają z ciągów
homotopii dla $O(q)\to O(q+1)\to S^q$.

Ustawiamy jednocześnie wszystkie rdzenie
w położeniu ogólnym, nie zmieniając ich klas.
Wymiar oczekiwanych przecięć własnych
i wzajemnych jest
$2k-(2k+1)=-1$; rdzenie są więc zanurzone
i parami rozłączne. Skończoność rodziny
pozwala zwęzić włókna $D^{k+1}$ tak, by
całe pogrubienia były rozłącznymi zanurzeniami
we wnętrzu $M$. Drogi bazowe wybieramy poza
rdzeniami: dla drogi i rdzenia wymiar przecięcia
to $1+k-(2k+1)<0$, a dla homotopii drogi
to $2+k-(2k+1)<0$, gdy $k\geq2$.

Każde $U_i=e_i(S^k\times D^{k+1})$ zwija się
na rdzeń. Najpierw cofamy $f|_{U_i}$ wzdłuż
tej retrakcji do $f\circ u_i$, a potem
używamy $H_i$ do homotopii z odwzorowaniem
stałym w punkcie bazowym $*$.
Ponieważ $U$ jest skończoną rozłączną sumą
zamkniętych otoczeń we wnętrzu $M$,
własność przedłużania homotopii pary
$(M,U\sqcup\partial M)$ przedłuża te
homotopie do $M$, bez zmiany mapy na brzegu.
Trywializację $F$ przenosimy wzdłuż
otrzymanej homotopii $f$.
\end{proof}

\begin{proof}
\emph{1. Przygotowanie i dwa lagranżjany.}
Chirurgie poniżej środka wymiaru z
\href{https://www.maths.ed.ac.uk/~v1ranick/books/scm.pdf}
{twierdzenia 1.2 Walla} przeprowadzamy we
wnętrzu $M$, aż $f$ stanie się $k$-spójne.
Wówczas
$G=\pi_{k+1}(f)\cong K_k(M;\Lambda)$
jest skończenie generowanym modułem.
Wybieramy generatory $x_1,\ldots,x_r$ wraz z
drogami bazowymi i zerującymi je homotopiami
w $X$. Lemat~\ref{lem:wall-odd-rozlaczne-sfery-22}
realizuje je przez rozłączne, obramowane
zanurzenia $S^k\times D^{k+1}$ we wnętrzu $M$.
Niech $U$ będzie sumą ich obrazów, a
$M_0=\overline{M\setminus\operatorname{int}U}$.
Po homotopii z tego samego lematu przyjmujemy
$f(U)=*$.
\href{https://www.maths.ed.ac.uk/~v1ranick/books/scm.pdf}
{Lemat 2.9 Walla} pozwala wybrać model $X$
z jedną najwyższą komórką, a przybliżenie
komórkowe przygotowuje odwzorowanie trójkątów
Poincarégo
\[
 (M;M_0,U)\longrightarrow(X;X_0,D^m).
\]
Ten krok jest potrzebny do stosowania dualności
do pary $(M_0,\partial U)$.

Grupa $H=K_k(\partial U;\Lambda)$ ma standardową
formę hiperboliczną $H_\epsilon(E)$, gdzie
$E\cong\Lambda^r$. Dwie strzałki brzegowe dają
\[
 L_0=\operatorname{im}\bigl(
 K_{k+1}(U,\partial U)\xrightarrow{\partial}H\bigr),
 \qquad
 L_1=\operatorname{im}\bigl(
 K_{k+1}(M_0,\partial U)\xrightarrow{\partial}H\bigr).
\]
Pierwszy moduł jest standardowym $E$. To, że
$L_1$ jest \emph{prostym lagranżjanem}, jest
konkretnym wejściem geometrycznym z dowodu
\href{https://www.maths.ed.ac.uk/~v1ranick/books/scm.pdf}
{lematu 5.7 Walla}, zastosowanego do
$(M_0,\partial U)$ po przygotowaniu trójkąta.
Prostą klasę baz i stabilną swobodność
$K_{k+1}(M_0,\partial U)$ zapewniają
\href{https://www.maths.ed.ac.uk/~v1ranick/books/scm.pdf}
{lemat 2.3(c) i wniosek z lematu 2.4 Walla}.
Dodanie generatorów zerowych zamienia moduł
stabilnie swobodny w swobodny z wybraną prostą
bazą. Wtedy wniosek 5.3.1 Walla przedłuża prosty
izomorfizm $L_0\to L_1$ do
$\alpha\in SU_r(\Lambda)$, przy czym
$\alpha(L_0)=L_1$.

Algebraicznie zawsze mamy dokładny ciąg
\[
 0\longrightarrow L_0\cap L_1
 \longrightarrow L_1
 \longrightarrow H/L_0
 \longrightarrow H/(L_0+L_1)
 \longrightarrow0.
\]
Jego środkowa strzałka wysyła $z$ na $[z]$;
jądro to $L_0\cap L_1$, a obraz to
$(L_0+L_1)/L_0$. Diagram Mayera--Vietorisa
dla $M=U\cup M_0$ identyfikuje dwa skrajne
moduły z $K_{k+1}(M)$ i $K_k(M)$. Po utożsamieniu
$H=E\oplus E^*$, $L_0=E$ i $L_1=\alpha(E)$
strzałkę środkową opisuje jedna macierz
\[
 T=\operatorname{pr}_{E^*}\circ\alpha|_E:
       E\longrightarrow E^*,
 \qquad K_{k+1}(M)\cong\ker T,
 \quad K_k(M)\cong\operatorname{coker}T.
\]
Są to jądra konkretnego wyboru $U$, a nie
jeszcze niezmiennik chirurgiczny.

\emph{2. Niezależność od wyborów i ruchy generatorów.}
Przy ustalonych zanurzeniach dwie dopuszczalne
proste bazy $L_1$ zmieniają $\alpha$ przez
mnożenie z prawej przez $TU_r$, ponieważ
oba automorfizmy przenoszą $E$ na $L_1$.
Regularna
homotopia zanurzeń $S^k\hookrightarrow M$
ma ślad będący immersją
$S^k\times I\to M\times I$.
Jej liczby przecięć wzajemnych i samoprzecięć
tworzą macierz $C=D-\epsilon D^*$, a
$\alpha$ zmienia się na $u_C\alpha$.
Odwrócenie tej konstrukcji przez lokalne ruchy
palcowe realizuje każde $u_C\in UU_r$.
Rachunek śladów, w tym jego geometryczną
realizację, przyjmujemy z
\href{https://www.maths.ed.ac.uk/~v1ranick/books/scm.pdf}
{§6 Walla, s. 60--61}.
Zatem przy jednej liście generatorów określony
jest dokładnie podwójny koset
$UU_r(\Lambda)\,\alpha\,TU_r(\Lambda)$.

Dwie skończone listy generatorów $G$ łączy
ciąg czterech ruchów. Można dołączyć drugą
listę do pierwszej, wypisać każdy jej element
jako kombinację pierwszych, a następnie
odwrócić rolę obu list. Każde dodanie
kombinacji rozkłada się na następujące kroki:
\begin{enumerate}
 \item[(T1)] Dołączamy albo usuwamy generator
 zerowy. Jego reprezentant mieści się w dysku
 rozłącznym z pozostałymi. Odpowiada temu
 $\alpha\mapsto\alpha\oplus\sigma$,
 jak pokazuje Lemat~\ref{lem:wall-odd-t1-22}.
 \item[(T2)] Permutujemy generatory. W bazie
 hiperbolicznej jest to sprzężenie przez
 $H(P)\in TU_r$, gdzie $P$ jest macierzą
 permutacji.
 \item[(T3)] Zastępujemy $x_1$ przez $u x_1$,
 gdzie $u=\pm g$ oraz $g\in\pi$. Zmiana
 orientacji i wybranej drogi daje sprzężenie
 przez $H(u)\in TU_r$; na odpowiadającej
 parze bazowej działa ono przez
 $e_1\mapsto u e_1$ i
 $f_1\mapsto(\bar u)^{-1}f_1
             =\pm w(g)g f_1$.
 \item[(T4)] Zastępujemy $x_1$ przez
 $x_1+x_2$. Łączymy dwie rurki łukiem
 reprezentującym właściwą klasę drogi,
 pogrubiamy łuk i przesuwamy uchwyt.
 Na module hiperbolicznym otrzymujemy
 \[
 e_1\mapsto e_1+e_2,\qquad e_2\mapsto e_2,
 \qquad f_1\mapsto f_1,\qquad f_2\mapsto f_2-f_1.
 \]
 Jest to $H(B)$ dla elementarnej macierzy
 $B\in\operatorname{Aut}^s_\Lambda(E)$.
 Geometryczna realizacja tego ruchu to
 konstrukcja uchwytowa cytowana przez Walla
 w §6 za jego wcześniejszą pracą [W8, s.~272].
\end{enumerate}

\begin{lemat}[Lokalny rachunek ruchu T1]
\label{lem:wall-odd-t1-22}
W sytuacji z części~1 dowodu niech $U$ będzie
sumą rurek reprezentujących wybraną listę
generatorów, a $\alpha\in SU_r(\Lambda)$
reprezentantem odpowiadającej pary lagranżjanów.
Dołączenie generatora zerowego przez standardowo
obramowaną sferę w dysku $D^{2k+1}\subset
\operatorname{int}(M\setminus U)$ pozwala wybrać
nowy reprezentant równy $\alpha\oplus\sigma$.
\end{lemat}
\begin{proof}
Na nieco większym dysku homotopujemy $f$ tak,
by na wybranym $D^{2k+1}$ było stałe, i przenosimy
przez tę homotopię dane normalne. Nowa sfera
z dyskiem ją ograniczającym przedstawia zero
w $\pi_{k+1}(f)$.
Lokalny model uzyskujemy z rozkładu
\[
 S^{2k+1}=\partial(D^{k+1}\times D^{k+1})
  =(S^k\times D^{k+1})\cup
    (D^{k+1}\times S^k)
\]
po wygładzeniu naroży i usunięciu małego dysku
z wnętrza drugiego składnika, rozłącznego z
$D^{k+1}\times\{*\}$. Pierwszy składnik
jest nową rurką $U_{\mathrm{new}}$, zawartą we
wnętrzu $D^{2k+1}$. Na jej brzegu
$S^k\times S^k$ wybieramy sfery
\[
 e=\{*\}\times S^k,\qquad
 f=S^k\times\{*\},
 \qquad\lambda(e,f)=1
\]
po dobraniu orientacji. Rdzeń nowej rurki
$S^k\times\{0\}$ łączymy walcem wewnątrz rurki
z $S^k\times\{*\}$; wraz z
$D^{k+1}\times\{*\}$ daje to, po wygładzeniu,
dysk ograniczający rdzeń.
Sfera $e$ jest brzegiem $\{*\}\times D^{k+1}$
w $U_{\mathrm{new}}$, a $f$ brzegiem
$D^{k+1}\times\{*\}$ w dopełnieniu tej rurki.
Obie mają produktowe obramowania, więc
$\mu(e)=\mu(f)=0$. Stąd lokalna forma ma bazę
hiperboliczną $(e,f)$, a obrazy strzałek
brzegowych z części~1 to odpowiednio
$L_{0,\mathrm{new}}=\Lambda e$ i
$L_{1,\mathrm{new}}=\Lambda f$.

Ponieważ $k\geq2$, sfera $\partial D^{2k+1}$
ma zerową homologię w stopniach $k$ i $k+1$.
Ciąg Mayera--Vietorisa po rozcięciu dopełnienia
wzdłuż tej sfery oddziela powyższy składnik
od starego diagramu definiującego $L_0,L_1$.
Dysk jest po prostu spójny; po wyborze drogi
bazowej wszystkie jego podniesienia dają ten
sam rachunek nad $\Lambda=\Z[\pi]$.
W otrzymanych identyfikacjach mamy zatem
\[
 H'=H\oplus(\Lambda e\oplus\Lambda f),\qquad
 L'_0=L_0\oplus\Lambda e,\qquad
 L'_1=L_1\oplus\Lambda f.
\]
Bazy pochodzące z obu produktowych dysków są
proste, a więc oba lagranżjany powiększają się
o prostą parę $(\Lambda e,\Lambda f)$.
Automorfizm przenoszący pierwszą prostą na drugą
możemy wybrać jako
$\sigma(e)=f$, $\sigma(f)=\epsilon e$.
Po dołączeniu tego wyboru do $\alpha$ otrzymujemy
$\alpha\oplus\sigma$.
Jest to lokalny rachunek ruchu (T1) z
\href{https://www.maths.ed.ac.uk/~v1ranick/books/scm.pdf}
{§6 Walla, s.~61--62}.
\end{proof}

\begin{lemat}[Rachunek zmian bazy w ruchach T2--T4]
\label{lem:wall-odd-t234-22}
Przy oznaczeniach z części~2 dowodu ruchy T2--T4
zmieniają bazę hiperboliczną przez $H(B)$,
gdzie $B\in\operatorname{Aut}^s_\Lambda(E)$ jest odpowiednio
permutacją, pomnożeniem jednego generatora przez
$u=\pm g$ albo elementarną zamianą
$e_1\mapsto e_1+e_2$. W ostatnim przypadku
baza dualna zmienia się według wzorów
$f_1\mapsto f_1$ i $f_2\mapsto f_2-f_1$.
\end{lemat}
\begin{proof}
Rozszczepienie~$H(B)$ z części~1 działa na parze
$(x,\varphi)\in E\oplus E^*$ wzorem
$H(B)(x,\varphi)=(Bx,\varphi\circ B^{-1})$.
Dlatego przy permutacji $e_i\mapsto e_{P(i)}$
funkcjonał $f_i$ przechodzi w $f_{P(i)}$;
pary hiperboliczne są po prostu przestawiane.

Jeśli $B(e_1)=u e_1$ i pozostałe $e_i$ są stałe,
to $B^{-1}(e_1)=u^{-1}e_1$. Z przyjętego
lewego działania na $E^*$ wynika
$H(B)(f_1)=(\bar u)^{-1}f_1$.
Dla $u=\pm g$ i inwolucji pierścienia grupowego
$\bar g=w(g)g^{-1}$ jest to
$(\bar u)^{-1}=\pm w(g)g$.
Jednocześnie $\lambda(u e_1,(\bar u)^{-1}f_1)=1$.

W ruchu T4 mamy
$B(e_1)=e_1+e_2$ i $B(e_2)=e_2$, zatem
$B^{-1}(e_1)=e_1-e_2$ i $B^{-1}(e_2)=e_2$.
Podstawienie do $f_i\circ B^{-1}$ daje
$H(B)(f_1)=f_1$ oraz $H(B)(f_2)=f_2-f_1$.
W szczególności
\[
 \lambda(e_1+e_2,f_2-f_1)=0,
 \qquad
 \lambda(e_1+e_2,f_1)=1,
 \qquad
 \lambda(e_2,f_2-f_1)=1.
\]
Macierze tych zmian są odpowiednio permutacyjne,
diagonalne z wpisem $\pm g$ i elementarne.
Ich klasy w $\operatorname{Wh}(\pi)$ znikają,
więc każde $B$ jest proste. Wzór na $H(B)$
pokazuje też zachowanie $\mu$, a tym samym
przynależność do $TU_r(\Lambda)$.
Geometryczne wykonanie ruchu T4 przez
dyfeomorfizm pogrubionego układu dwóch rurek
pozostaje wejściem z
\href{https://www.maths.ed.ac.uk/~v1ranick/books/scm.pdf}
{§6 Walla, s.~62--63, odsyłającym do [W8, s.~272]}.
\end{proof}

Ruch $T4$ realizuje macierze elementarne,
$T3$ mnożenie generatora przez $\pm g$,
a $T2$ permutacje. Z definicji zerowej
klasy w $\operatorname{Wh}(\pi)$ każda prosta
zmiana bazy po stabilizacji rozkłada się
na takie operacje; nie trzeba zakładać,
że sama jest iloczynem macierzy elementarnych.
Ich wpływ znika w ilorazie przez $TU$.
Ruch $T1$ znika w
ilorazie przez $RU$, gdyż $\sigma\in RU$.
To wyjaśnia niezależność klasy
$[\alpha]\in SU/RU$ od generatorów i
reprezentujących je zanurzeń. W szczególności
stabilizacja geometryczna dodaje $\sigma$,
a nie identyczność; dopiero po przejściu
do $SU/RU$ obie dają ten sam wynik.

\emph{3. Chirurgia, kobordyzm i kryterium zera.}
\begin{lemat}[Lokalny efekt chirurgii środkowej]
\label{lem:wall-odd-surgery-22}
Niech $k\geq2$ i niech $U_1=S^k\times D^{k+1}$
będzie pierwszą z obramowanych rurek w $U$.
Po chirurgii na jej rdzeniu można wybrać bazę
hiperboliczną nowego $H=K_k(\partial U)$ tak,
by reprezentant pary lagranżjanów zmienił się
z $\alpha$ na $\sigma_1\alpha$, gdzie
$\sigma_1=\sigma\oplus 1$ działa na pierwszej
płaszczyźnie hiperbolicznej. Pozostałe składniki
$U$ i lagranżjan dopełnienia $L_1$ nie zmieniają się.
\end{lemat}
\begin{proof}
Chirurgia usuwa $S^k\times\operatorname{int}D^{k+1}$
i wkleja $U'_1=D^{k+1}\times S^k$ wzdłuż tego samego
brzegu $S^k\times S^k$. Oznaczając
$U' = U'_1\sqcup U_2\sqcup\cdots\sqcup U_r$,
otrzymujemy $M'=M_0\cup_{\partial U}U'$.
Dopełnienie $M_0$ i jego strzałka brzegowa
\[
 K_{k+1}(M_0,\partial U)\xrightarrow{\partial}
 K_k(\partial U)
\]
pozostają takie same. Zmienia się strzałka
od strony rurki: dysk $\{*\}\times D^{k+1}$
w $U_1$ miał brzeg $e_1$, natomiast dysk
$D^{k+1}\times\{*\}$ w $U'_1$ ma brzeg $f_1$.
Zatem w dawnych współrzędnych
\[
 L'_0=\Lambda f_1\oplus
       \bigoplus_{i=2}^r\Lambda e_i,
 \qquad L'_1=L_1=\alpha(E).
\]
Wybierzmy identyfikację nowych współrzędnych
ze starymi daną przez $j=\sigma_1^{-1}$.
Przenosi ona nowe standardowe $E$ na $L'_0$;
znak generatora $f_1$ wybieramy zgodnie z $j$.
W nowych współrzędnych drugi lagranżjan to
$j^{-1}L'_1=\sigma_1\alpha(E)$, skąd można
wybrać $\alpha'=\sigma_1\alpha$.
Gdy $\epsilon=-1$, wybór identyfikacji
$j=\sigma_1$ zamiast $\sigma_1^{-1}$ zmienia
wynik tylko o $H(-1)\in TU_r(\Lambda)$
na pierwszej parze; gdy $\epsilon=1$, obie
identyfikacje są równe. Jest to rachunek z
\href{https://www.maths.ed.ac.uk/~v1ranick/books/scm.pdf}
{dowodu lematu 6.1 Walla, s.~63}; Wall działa
macierzami na bazach z prawej, więc zapisuje
odpowiadający czynnik po prawej stronie.
\end{proof}

Dokładne stwierdzenie obejmujące także ruchy
odwrotne podaje
\href{https://www.maths.ed.ac.uk/~v1ranick/books/scm.pdf}
{lemat 6.1 Walla}: chirurgia do prostej
równoważności jest możliwa wtedy i tylko wtedy,
gdy $\alpha$ jest dwustronnie równoważne
bezpośredniej sumie $\Sigma$ kopii $\sigma$
przy użyciu $RU$.

Wyjaśnijmy, dlaczego lemat obejmuje dowolny
kobordyzm normalny między przygotowanymi
$k$-spójnymi mapami. Dany kobordyzm
$(N;M,M')\to X\times I$ przygotowuje się
chirurgiami tak, aby odwzorowanie było
$(k+1)$-spójne. Jego moduł
$K_{k+1}(N,M)$ jest stabilnie swobodny i ma
prostą bazę, na mocy argumentu z lematu 2.3
Walla. Chirurgia na trywialnej $k$-sferze
we wnętrzu $N$ dodaje do tego modułu swobodny
składnik rangi dwa; po skończonej liczbie takich
kroków moduł ma rzeczywistą prostą bazę.
Twierdzenie o rozkładzie na uchwyty cytowane
przez Walla jako [W13, IV, 5.5] pozwala
reprezentować tę bazę uchwytami wymiaru
$k+1$ dołączonymi do $M$. Pozostała część
jego rozkładu jest prostym $h$-kobordyzmem,
czyli $s$-kobordyzmem; stosujemy
Twierdzenie~\ref{thm:s-kobordyzm-20}.
W efekcie każda zmiana przez kobordyzm
normalny rozkłada się na opisane chirurgie
środkowe, zmiany wyborów oraz składnik
iloczynowy. Jest to geometryczna część
dowodu lematu 6.1 Walla, użyta tu jako
wynik zewnętrzny.

Normalność $RU$ wynika z cytowanych już
lematu 6.2 i twierdzenia 6.3 Walla. Możemy
więc określić
\[
 \theta(f,F)=[\alpha]\in SU(\Lambda)/RU(\Lambda).
\]
Poprzednie dwa akapity dowodzą jego
niezależności od przygotowania i niezmienniczości
względem kobordyzmu normalnego relatywnie do
brzegu. Jeśli $\theta=0$, algebraiczny krok
lematu 6.1, korzystający z rozkładu elementów
$RU$ na $\sigma$, $UU$ oraz hiperboliczne
automorfizmy elementarnych zmian bazy,
zamienia $\alpha$ za pomocą dozwolonych
ruchów w $\Sigma$. Dla $\alpha=\Sigma$ oba
lagranżjany są dopełniające, a strzałka
$T:E\to E^*$ jest \emph{prostym} izomorfizmem.
Z powyższego ciągu wynikają
$K_{k+1}(M)=K_k(M)=0$. Niższe jądra
zniknęły podczas przygotowania, a wyższe
znikają przez dualność Poincarégo.
\href{https://www.maths.ed.ac.uk/~v1ranick/books/scm.pdf}
{Wniosek z lematu 2.4 Walla} kontroluje
prostotę powstałego acyklicznego kompleksu
jądra, toteż końcowe odwzorowanie jest
prostą równoważnością homotopijną.
W odwrotną stronę dla prostej równoważności
wybieramy $U=\varnothing$: macierz $\alpha$
ma rozmiar zero i jej klasa jest zerowa.

Jeżeli dane przedłużymy przez prostą
równoważność na dołączonej parze Poincarégo,
to samo $U$ oblicza przeszkodę po zmianie
współczynników
$\Z[\pi_1(X)]\to\Z[\pi_1(X')]$.
Macierz $\alpha$ przechodzi w macierz otrzymaną
przez zastosowanie tego homomorfizmu do
każdego wyrazu. To daje naturalność.
W przypadku brzegu wszystkie zanurzenia i
chirurgie wykonujemy poza ustalonym kołnierzem
$\partial M\times[0,1]$, a $s$-kobordyzm
stosujemy relatywnie do tego kołnierza.
\href{https://www.maths.ed.ac.uk/~v1ranick/books/scm.pdf}
{Wall, §6, s.~58 i twierdzenie 6.4}
sformułował wynik relatywny, choć sam rachunek
§6 prowadzi bez zapisywania brzegu.
\end{proof}

Opis grupy nieparzystej przez klasy formacji
$(H;L_0,L_1)$ jest równoważny powyższemu
opisowi macierzowemu $SU/RU$. Równoważność
ta korzysta z relacji formacji brzegowych;
jest osobnym wynikiem
\href{https://www.maths.ed.ac.uk/~v1ranick/books/scm.pdf}
{Ranickiego przytoczonym przez Walla po
twierdzeniu 6.6, w uwagach „Forms and formations”}.
Wcześniejszy przykład z $T=2$ można więc
odczytywać zarówno jako formację brzegową,
jak i jako zerową klasę w $SU/RU$.

\begin{uwaga}[Dlaczego wymiar pięć jest objęty kryterium]
\label{uwaga:wall-odd-five-22}
Dla $m=5$ mamy $k=2$. Ogólne położenie
rozdziela skończenie wiele reprezentujących
generatorów sfer $S^2$ w $M^5$, ponieważ
$2+2<5$. Normalna wiązka takiej sfery ma
rangę trzy. Stabilne obramowanie pochodzące
z danych normalnych daje rzeczywiste
obramowanie, gdyż
$\pi_1(SO(3))\to\pi_1(SO)$ jest izomorfizmem;
stąd argument z osadzaniem obramowanych
sfer działa już przy $k=2$. Ślad chirurgii
ma wymiar sześć, więc stosowane w dowodzie
gładkie twierdzenie o $s$-kobordyzmie jest
w swoim zakresie wymiarów. W szczególności
twierdzenie 6.4 Walla obejmuje $m=5$.
Osobne \emph{twierdzenie realizacyjne} 6.5
Walla zakłada $k\geq3$, czyli $m\geq7$:
nie wolno jego zakresu przenosić na
powyższe kryterium zaniku przeszkody.
\end{uwaga}

\begin{uwaga}[Gdzie kończy się argument z homologii]
\label{uwaga:homologia-L-22}
Zerowa zwykła homologia jądra nie wystarcza:
nad $\Z[\pi]$ mogą pozostać nietrywialne
klasy dróg. Nie wystarcza również nieosobliwość
formy: potrzeba lagranżjanu zgodnego z $\mu$.
Wreszcie istnienie \emph{homotopijnej}
równoważności nie rozstrzyga automatycznie
kwestii torsji Whiteheada potrzebnej dla
równoważności prostej.
\end{uwaga}

\section{Trzy rachunki i kolejny krok}
\label{sec:przyklady-L-22}

\begin{enumerate}[label=(\arabic*)]
\item Dla identyczności $\id:X\to X$ jądro
kompleksu mapy jest ściągalne, więc
$\theta(\id)=0$.
\item Dla normalnej mapy z
Przykładu~\ref{prz:produkt-L-22} środkowe
jądro ma rangę $2$, lecz jest hiperboliczne.
Dlatego jego klasa $L$ znika.
\item Algebraiczne formy $E_8$ oraz para
antysymetryczna z $\mu(e)=\mu(f)=1$
pokazują odpowiednio niezerową przeszkodę
sygnaturową i niezerową przeszkodę Arfa.
Jeśli taka forma jest jądrem mapy normalnej,
jej niezerowa klasa wyklucza końcową
równoważność w tej normalnej klasie kobordyzmu.
\end{enumerate}

Rozdział~\ref{ch:chirurgia-geometryczna-21}
dał \emph{operację}, a ten rozdział daje
\emph{przeszkodę}. Do klasyfikowania
rozmaitości pozostaje zestawić normalne
niezmienniki, klasy struktur i grupy $L$
w ciągu dokładnym chirurgii. W takim ciągu
mapa do grupy $L$ mówi, które dane normalne
można zrealizować przez równoważność
homotopijną, a działanie poprzedniej grupy
$L$ opisuje niejednoznaczność wyniku.

\par\smallskip
{\footnotesize\noindent\emph{Dalsza lektura:}
\href{https://webhomes.maths.ed.ac.uk/~v1ranick/books/scm.pdf}
{C.~T.~C. Wall, \emph{Surgery on Compact Manifolds},
rozdziały o formach i przeszkodach},
\href{https://webhomes.maths.ed.ac.uk/~v1ranick/papers/ats.pdf}
{A.\ Ranicki, \emph{The Algebraic Theory of Surgery},
§7--8},
\href{https://link.springer.com/book/10.1007/978-3-031-56334-8}
{W.\ Lück i T.\ Macko, \emph{Surgery Theory:
Foundations}, rozdziały o teorii $L$}.\par}

\clearpage
